\documentclass[11pt]{article}
\usepackage[a4paper,margin=27mm]{geometry}
\usepackage{amsmath,amssymb,amsthm,mathtools}
\usepackage{booktabs}
\usepackage{array}
\usepackage{graphicx}
\usepackage{float}
\usepackage{microtype}
\usepackage{xcolor}
\usepackage[hidelinks]{hyperref}
\usepackage{enumitem}

\newtheorem{theorem}{Theorem}[section]
\newtheorem{proposition}[theorem]{Proposition}
\newtheorem{corollary}[theorem]{Corollary}
\newtheorem{lemma}[theorem]{Lemma}
\newtheorem{remark}[theorem]{Remark}
\newcommand{\E}{\mathbb E}
\newcommand{\R}{\mathbb R}
\newcommand{\I}{\mathrm I}
\newcommand{\tr}{\operatorname{tr}}
\newcommand{\diag}{\operatorname{diag}}
\newcommand{\sech}{\operatorname{sech}}
\hypersetup{
 pdftitle={The Schwarzian Bridge in a Single Sigmoid Neuron: Oracle Gram Resolution and Spherical Saturation Phases},
 pdfauthor={Lluis Eriksson},
 pdfkeywords={single neuron, sigmoid, Loewner matrix, matrix monotonicity, Fisher information, saturation}
}

\title{\bfseries The Schwarzian Bridge in a Single Sigmoid Neuron\\
\large Oracle Gram Resolution and Spherical Saturation Phases}
\author{Lluis Eriksson\\Independent Researcher\\\texttt{lluiseriksson@gmail.com}}
\date{August 2026}

\begin{document}
\maketitle

\begin{abstract}
A sigmoid is increasing on scalars, yet its bend controls two different geometries.  The classical Schwarzian
criterion implies that a negative-Schwarzian activation is not matrix-monotone of order two.  We give a closed
logistic formula showing every distinct two-point L\"owner determinant to be negative, a dimension-minimal
$2\times2$ positive-semidefinite perturbation witness certified with interval arithmetic, and the corresponding
classical serial-chain consequence.

Our central link is that the same local invariant determines the leading quadratic onset coefficient of saturation
anisotropy.  At a centered inflection of an increasing activation $g$, for a sensitivity $h_p=(g')^p$ and
isotropic spherically symmetric input $X$,
\[
 \kappa_{p,X}(r)=1-q_Xp\,Sg(0)r^2+O(r^4),\qquad
 q_X=\frac{\E\lVert X\rVert^4}{d(d+2)}.
\]
The Gaussian coefficient is the special case $q_X=1$.  The adjacent-point L\"owner determinant is
$g'(0)^2Sg(0)\delta^2/6+O(\delta^3)$.  Negative Schwarzian therefore simultaneously gives local
noncommutative order failure and positive initial radial--tangential anisotropy.  We turn the equality of limits
into a finite-scale inequality with explicit constants determined by derivative suprema and input moments.

For isotropic Gaussian inputs in dimension $d\ge2$, population curvature at a teacher of norm $r$ has one radial
eigenvalue and $d-1$ equal tangential eigenvalues.  For any nonnegative sensitivity profile $h$ with finite defining expectations, these
eigenvalues are exactly
\[
 \alpha_h(r)=\E h(rZ),\qquad \beta_h(r)=\E[Z^2h(rZ)],\qquad Z\sim N(0,1),
\]
and, if its zeroth and second Lebesgue moments are finite and positive,
\[
 \alpha_h(r)\sim \frac{\varphi(0)m_0}{r},\qquad
 \beta_h(r)\sim \frac{\varphi(0)m_2}{r^3}.
\]
Thus saturation suppresses magnitude information by two additional powers: the radial--tangential ratio obeys
$\alpha_h(r)/\beta_h(r)\sim(m_0/m_2)r^2$.  For every real $p>0$ and
$h_p=\sigma'^p$, we prove strict monotonicity and the unified identities
\[
 m_0(p)=\frac{\Gamma(p)^2}{\Gamma(2p)},\qquad
 \kappa_p(r)\sim\frac{r^2}{2\psi_1(p)},\qquad
 \kappa_p(r)=1+\frac p2r^2-\frac p8r^4+\frac{p(p+1)}{16}r^6+O(r^8).
\]
For $n$ samples we prove an explicit two-sided relative L\"owner bound for the fixed-teacher oracle weighted
Gram matrix.  A conditional-Gaussian lower bound shows that the sample complexity of this specific loss is sharply
\[
 \Theta_p\!\left(\frac{r[d+\log(1/\delta)]}{\varepsilon^2}\right).
\]
The factor $rd$ is intrinsic: a tangential off-diagonal column has $d-2$ coordinates and normalized variance of
order $r/n$.  We also give a uniform full-dyadic-shell extension and concrete sufficient constants
$C_1=22929$ and $C_2=294162$.  The weaker bottom-eigenspace target has a different extreme-saturation law:
its empirical projector converges to the normal of the $d-1$ samples nearest the transition hyperplane, with exact
angle distribution $\tan\theta_\infty\stackrel d=\lVert\zeta\rVert/|G|$, matching
$\Theta(d^{3/2}/\varepsilon)$ complexity at fixed nondegenerate confidence, and matching $1/\delta$ dependence
in the theorem's stated confidence range.

The quadratic saturation law extends beyond Gaussian inputs.  If $X=RU$ is isotropic spherical,
$\E R^{-1}<\infty$, then
\[
 \kappa_h(r)\sim Q_R\frac{m_0}{m_2}r^2,\qquad
 Q_R=\frac{\E R}{(d-1)\E R^{-1}}.
\]
Gaussian, fixed-sphere and isotropic Student inputs give respectively $Q_R=1$, $d/(d-1)$ and
$(\nu-2)/(\nu-1)$.  For radii with density $f_R(\rho)\sim c\rho^{a-1}$ at zero, we prove exact leading constants
in all three phases: $r^{-(a+2)}$ for $0<a<1$, $r^{-3}\log r$ at $a=1$, and $r^{-3}$ for $a>1$.
An autonomous Abelian kernel lemma handles the critical boundary, and deterministic quadrature exhibits all
three phases.  We also give an isotropic Rademacher counterexample.  These claims refine, but do not replace, prior one-sided and
bilateral Hessian results and recent labeled-estimation theory.  Deterministic replays, a fixed-seed finite-sample
phase diagram and an immutable public source record accompany the proofs.
\end{abstract}

\paragraph{Keywords.} single neuron; sigmoid; matrix monotonicity; L\"owner order; Schwarzian derivative;
Fisher information; spherical design; weighted sample covariance; finite-sample concentration.

\section{Question, result, and claim boundary}

The scalar logistic curve is
\[
 \sigma(s)=\frac1{1+e^{-s}},\qquad
 \sigma'(s)=\frac1{4}\sech^2(s/2).
\]
It is strictly increasing, so larger scalar scores produce larger outputs.  We first ask whether this order survives
spectral functional calculus on noncommuting matrices.  It does not, at the smallest possible matrix size and on
every interval.  We then return to the ordinary teacher neuron $\sigma(w_\star^\top X)$ with
$X\sim N(0,I_d)$.  There, a perturbation perpendicular to $w_\star$ rotates the separating hyperplane, while a
parallel perturbation changes its confidence scale and must be detected through samples near the transition layer.
Our goal is to compute both failures exactly.

The main contributions are:
\begin{enumerate}[leftmargin=2em]
 \item a closed logistic specialization of the classical negative-Schwarzian obstruction, proving directly that
 every nontrivial two-point L\"owner matrix is indefinite;
 \item a dimension-minimal $2\times2$ witness $A\preceq B$ but $\sigma(A)\npreceq\sigma(B)$, with
 outward-rounded eigenvalue certification, and a L\"owner-order reformulation of the known serial-composition
 consequence;
 \item a spherical-input Schwarzian bridge equating the local order defect with the leading
 radial--tangential anisotropy, including its exact fourth-moment multiplier $q_X$;
 \item an explicit finite-scale bridge bounding the discrepancy between the normalized L\"owner determinant
 and anisotropy at nonzero $(\delta,r)$;
 \item a profile-generic two-eigenvalue refinement under Gaussian design, with finite-$r$ brackets and, for
 every $p>0$, closed moments, strict ratio monotonicity and exact endpoint laws for $h_p=\sigma'^p$;
 \item a two-sided relative concentration theorem for the fixed-teacher oracle sensitivity matrix, with a matching
 $r[d+\log(1/\delta)]/\varepsilon^2$ lower bound for bilateral relative Loewner loss, explicit constants for
 $p=1,2$, an autonomous high-confidence lower-tail lemma, and a uniform full-dyadic-shell extension;
 \item a spectral-lexicography theorem for the weaker bottom-eigenspace target, including an exact limiting angle
 law, matching $\Theta(d^{3/2}/\varepsilon)$ complexity at fixed nondegenerate confidence, and an explicit
 high-confidence range in the iterated extreme-saturation experiment;
 \item a spherical saturation theorem with exact asymptotic radius factor $Q_R$, an elliptical whitening
 corollary, and exact constants and a critical logarithmic law for radii regularly varying at zero; and
 \item exact local consequences for stationary fixed-step optimization and estimation, plus a documented replay sequence
 of the identities,
 interval witness, inequalities, and figure.
\end{enumerate}

The matrix-order theorem concerns spectral functional calculus $\sigma(A)$, not entrywise activation of a vector.
The radial theorem is not a claim about deep networks, global optimization, implicit bias, or arbitrary inputs.  Its
local coefficient and its large-saturation law extend from Gaussian to spherical inputs under different, stated
moment hypotheses.  The
two questions meet at a precise boundary: scalar monotonicity controls neither noncommuting PSD perturbations nor
the conditioning of parameter directions.

\section{A sigmoid is not an order-preserving operator}

For Hermitian matrices, $A\preceq B$ means that $B-A$ is positive semidefinite.  A scalar function $f$ is
\emph{matrix-monotone of order two} on an interval $J$ if $f(A)\preceq f(B)$ for every pair of $2\times2$
Hermitian matrices with spectra in $J$ and $A\preceq B$, where $f(A)$ is defined by spectral functional calculus.
If $A$ and $B$ commute, scalar monotonicity is sufficient because they are simultaneously diagonalizable.  The
question is what happens without commutativity.

For distinct $x,y\in J$, define the two-point L\"owner matrix
\begin{equation}
 L_f(x,y)=
 \begin{pmatrix}
  f'(x)&\dfrac{f(x)-f(y)}{x-y}\\[5pt]
  \dfrac{f(x)-f(y)}{x-y}&f'(y)
 \end{pmatrix}.
 \label{eq:loewnermatrix}
\end{equation}
The fixed-order L\"owner criterion says that matrix monotonicity of order two requires
$L_f(x,y)\succeq0$ for every pair \cite{hiai2010,heinavaara2019}.  The equivalent local criterion
$Sf\ge0$ for increasing smooth functions is classical \cite[Lemma 10]{kozlovski2008}; the result below is an
explicit logistic specialization, not a new qualitative matrix-monotonicity criterion.

Cook, Hammerlindl and Tucker define
\[
 \chi_f(x,y)=f'(x)f'(y)
 \left(\frac{x-y}{f(x)-f(y)}\right)^2
\]
and study the class $\chi_f\le1$, including sigmoid activations and serial compositions \cite{cook2023}.  If
$D_f(x,y)=(f(x)-f(y))/(x-y)$, then
\begin{equation}
 \det L_f(x,y)=D_f(x,y)^2\,[\chi_f(x,y)-1].
 \label{eq:coexpanding}
\end{equation}
Thus their two-point inequality is algebraically the nonpositive L\"owner-determinant inequality.  Our contribution
in this subsection is the strict closed logistic identity and the certified matrix witness, not the underlying
Schwarzian/composition mechanism.

\begin{theorem}[closed logistic two-point determinant]
\label{thm:loewner}
Let
\[
 s_{\alpha,b}(x)=\frac1{1+e^{-\alpha(x-b)}},\qquad \alpha>0.
\]
For every $x\ne y$, the L\"owner matrix $L_{s_{\alpha,b}}(x,y)$ has negative determinant.  Consequently the
logistic activation is not matrix-monotone of order two on any nondegenerate interval.
\end{theorem}

\begin{proof}
Put $u=e^{\alpha(x-b)}$, $v=e^{\alpha(y-b)}$, and $r=v/u=e^{\alpha(y-x)}$.  Direct cancellation gives the exact
identity
\begin{equation}
 \det L_{s_{\alpha,b}}(x,y)
 =\frac{\alpha^2u^2}{(1+u)^2(1+v)^2}
 \left[r-\frac{(r-1)^2}{(\log r)^2}\right].
 \label{eq:loewnerdet}
\end{equation}
For $r\ne1$, write $t=\log r$.  Strict convexity of $\sinh$ on $(0,\infty)$ gives
\[
 |r-1|=2e^{t/2}|\sinh(t/2)|>e^{t/2}|t|=\sqrt r\,|\log r|.
\]
The bracket in \eqref{eq:loewnerdet} is therefore strictly negative.  Every nondegenerate interval contains a
distinct pair, so the L\"owner criterion completes the proof.
\end{proof}

This is stronger than finding one bad scale: changing the slope $\alpha$, shifting the bias $b$, or restricting the
score range never removes the obstruction.  It also identifies the precise failure.  The diagonal entries in
\eqref{eq:loewnermatrix} are positive, but the divided difference is too large for their geometric mean.

\begin{proposition}[dimension-minimal explicit PSD witness]
\label{prop:witness}
For the standard logistic sigmoid, let
\begin{equation}
 A=\begin{pmatrix}-1/2&0\\0&1/2\end{pmatrix},\qquad
 B=A+\frac1{100}\begin{pmatrix}1&1\\1&1\end{pmatrix}.
 \label{eq:witness}
\end{equation}
Then $A\preceq B$ but $\sigma(A)\npreceq\sigma(B)$.  At 256-bit outward-rounded Arb precision,
\begin{align}
 \operatorname{spec}(B)&\subset[-0.4901000,-0.4900999]\cup[0.5100999,0.5101000],\\
 \operatorname{spec}(\sigma(B)-\sigma(A))
 &\subset[-9.9150,-9.9147]\,10^{-5}\ \cup\ [0.0047990,0.0047992],\label{eq:witnessspec}\\
 \det(\sigma(B)-\sigma(A))&\in[-4.7583,-4.7581]\,10^{-7}.
\end{align}
\end{proposition}

\begin{proof}
$B-A$ has eigenvalues $0$ and $1/50$.  If $a=1/2$, $\varepsilon=1/100$, and
$d=(a^2+\varepsilon^2)^{1/2}$, then $B$ has eigenvalues $\lambda_\pm=\varepsilon\pm d$ and
\[
 \sigma(B)=mI+n(B-\varepsilon I),\quad
 m=\frac{\sigma(\lambda_+)+\sigma(\lambda_-)}2,\quad
 n=\frac{\sigma(\lambda_+)-\sigma(\lambda_-)}{2d}.
\]
Substitution reduces the three enclosures to scalar exponential and square-root balls.  The replay evaluates them
with directed Arb rounding; the negative eigenvalue in \eqref{eq:witnessspec} proves the claim.
\end{proof}

The witness is illustrative, while Theorem~\ref{thm:loewner} is the proof for every interval.  The derivative form
also explains why an indefinite L\"owner matrix produces witnesses: for diagonal
$A=\diag(x,y)$, the Fr\'echet derivative in the all-ones PSD direction is precisely the Schur product with
$L_f(x,y)$.

\subsection{Serial depth does not repair the order}

For a $C^3$ function with nonzero derivative, its Schwarzian is
\begin{equation}
 Sf=\frac{f'''}{f'}-\frac32\left(\frac{f''}{f'}\right)^2.
 \label{eq:schwarzian}
\end{equation}
The logistic sigmoid and a scaled hyperbolic tangent satisfy
\begin{equation}
 S s_{\alpha,b}=-\frac{\alpha^2}{2},\qquad
 S\tanh(\alpha x+\beta)=-2\alpha^2.
 \label{eq:schwarzianvalues}
\end{equation}

\begin{corollary}[serial spectral reformulation]
\label{thm:serial}
Let $F$ be a finite scalar composition of positive-slope affine maps and at least one increasing logistic or
hyperbolic-tangent activation.  Then $SF<0$ everywhere, and $F$ is not matrix-monotone of order two on any
nondegenerate interval.
\end{corollary}

\begin{proof}
The composition identity
\begin{equation}
 S(f\circ g)=(Sf\circ g)(g')^2+Sg
 \label{eq:schwarzchain}
\end{equation}
and \eqref{eq:schwarzianvalues} give $SF<0$, since positive affine maps have zero Schwarzian and all derivatives in
the chain are positive.  For $h\to0$, Taylor expansion of the two-point determinant gives
\begin{equation}
 \det L_F(x,x+h)=\frac{F'(x)^2}{6}\,SF(x)\,h^2+O(h^3).
 \label{eq:localschwarz}
\end{equation}
It is negative for all sufficiently small nonzero $h$.  Every interval contains such a pair.
\end{proof}

The corollary restates, in spectral L\"owner order, the serial closure proved in the nowhere-coexpanding framework
\cite{cook2023}.  It is deliberately about \emph{serial scalar chains}.  It does not cover sums of neurons, skip
connections, multivariate architectures, or entrywise coordinate order.  Its content is that composing more of the
same S-shaped bend cannot turn the resulting scalar function into a PSD-order-preserving spectral activation.

\section{Two models, one sensitivity matrix}

Write $r=\lVert w_\star\rVert$ and $u=w_\star/r$ when $r>0$.  We use two standard observation models.

\paragraph{Squared-output regression.}
For the realizable population loss
\begin{equation}
 \mathcal L_{\rm sq}(w)=\frac12\E\bigl[\sigma(w^\top X)-\sigma(w_\star^\top X)\bigr]^2,
 \label{eq:squareloss}
\end{equation}
the residual vanishes at $w_\star$, so its Hessian is the Gauss--Newton matrix
\begin{equation}
 H_{\rm sq}(w_\star)=\E\bigl[\sigma'(w_\star^\top X)^2XX^\top\bigr].
 \label{eq:squarehess}
\end{equation}
The same matrix, divided by a known output-noise variance $\tau^2$, is the Fisher information of
$Y=\sigma(w_\star^\top X)+\varepsilon$, $\varepsilon\sim N(0,\tau^2)$.

\paragraph{Bernoulli logistic neuron.}
If $Y\mid X\sim\mathrm{Bernoulli}(\sigma(w_\star^\top X))$, the expected negative-log-likelihood Hessian and the
Fisher information coincide:
\begin{equation}
 H_{\rm B}(w_\star)=\E\bigl[\sigma'(w_\star^\top X)XX^\top\bigr].
 \label{eq:bernhess}
\end{equation}

Both are instances of the sensitivity matrix
\begin{equation}
 H_h(r,u):=\E\bigl[h(r u^\top X)XX^\top\bigr], \qquad h\ge0.
 \label{eq:profilematrix}
\end{equation}
The power of the reduction below is that no high-dimensional integration remains.

Amari, Karakida and Oizumi already derived the exact one-unit Gaussian Fisher decomposition for
$h=\phi'^2$, including the radial--tangential split and the bias block \cite{amari2019}.  We do not claim that
decomposition or the squared-output branch as new.
For the Bernoulli profile $h=\sigma'$, Chen and Mazumdar recently identified the same radial and orthogonal
Hessian functions and proved their $r^{-3}$ and $r^{-1}$ orders in a finite-sample analysis of logistic regression
\cite{chen2026gd}.  Their follow-up proves the minimax norm-estimation rate $\Theta(\sqrt{r^3/n})$
\cite{chen2026minimax}.  We do not claim those exponents or the Bernoulli optimization penalty as new.  The
purpose of the next sections is to isolate the profile-generic mechanism, give finite-radius brackets and exact
logistic constants, prove strict monotonicity of both eigenvalue ratios, and connect their small-$r$ coefficients to
the Schwarzian order defect.

\section{Universal Gaussian saturation theorem}

Let $\varphi(z)=(2\pi)^{-1/2}e^{-z^2/2}$ and define the even sensitivity moments
\[
 m_j(h):=\int_{\R} t^j h(t)\,dt\quad(j=0,2,4),
\]
whenever finite.  Evenness of $h$ is not needed for the spectral reduction or leading asymptotics, but it holds in
our sigmoid applications and makes the interpretation symmetric.

\begin{theorem}[exact spectrum and finite-$r$ brackets]
\label{thm:universal}
Let $d\ge2$, $X\sim N(0,I_d)$, $r>0$, $u\in S^{d-1}$, and let $h:\R\to[0,\infty)$ be measurable with
$\E[(1+Z^2)h(rZ)]<\infty$.  Then
\begin{equation}
 H_h(r,u)=\alpha_h(r)(I-uu^\top)+\beta_h(r)uu^\top,
 \label{eq:decomp}
\end{equation}
where
\begin{equation}
 \alpha_h(r)=\E h(rZ),\qquad \beta_h(r)=\E[Z^2h(rZ)],\qquad Z\sim N(0,1).
 \label{eq:alphabeta}
\end{equation}
If $m_0,m_2<\infty$ and are positive, then
\begin{equation}
 \alpha_h(r)\sim\frac{\varphi(0)m_0}{r},\qquad
 \beta_h(r)\sim\frac{\varphi(0)m_2}{r^3},\qquad
 \kappa_h(r):=\frac{\alpha_h(r)}{\beta_h(r)}\sim\frac{m_0}{m_2}r^2.
 \label{eq:univasymp}
\end{equation}
If $m_0,m_2,m_4<\infty$, the following nonasymptotic brackets hold for every $r>0$:
\begin{align}
 \frac{\varphi(0)}r\left(m_0-\frac{m_2}{2r^2}\right)
 &\le \alpha_h(r)\le \frac{\varphi(0)m_0}{r},\label{eq:alphabounds}\\
 \frac{\varphi(0)}{r^3}\left(m_2-\frac{m_4}{2r^2}\right)
 &\le \beta_h(r)\le \frac{\varphi(0)m_2}{r^3}.
 \label{eq:betabounds}
\end{align}
Negative lower endpoints are interpreted as valid but uninformative bounds.
\end{theorem}

\begin{proof}
Decompose $X=Zu+Y$, where $Z\sim N(0,1)$, $Y\sim N(0,I-uu^\top)$, and $Z,Y$ are independent.
The mixed terms vanish and $\E YY^\top=I-uu^\top$, giving \eqref{eq:decomp}--\eqref{eq:alphabeta}.
Changing variables $t=rz$ gives the exact one-dimensional formulae
\begin{equation}
 \alpha_h(r)=\frac{\varphi(0)}r\int_\R h(t)e^{-t^2/(2r^2)}dt,
 \quad
 \beta_h(r)=\frac{\varphi(0)}{r^3}\int_\R t^2h(t)e^{-t^2/(2r^2)}dt.
 \label{eq:scaledintegrals}
\end{equation}
Dominated convergence proves \eqref{eq:univasymp}.  Finally,
$1-x\le e^{-x}\le1$ for $x\ge0$ inserted into \eqref{eq:scaledintegrals} proves
\eqref{eq:alphabounds}--\eqref{eq:betabounds}.
\end{proof}

\section{Spherical saturation beyond Gaussian design}
\label{sec:global-spherical}

The Gaussian density at the transition hyperplane is not the essential constant.  Let $X=RU$ be spherically
symmetric, where $U$ is uniform on $S^{d-1}$ and independent of the radius $R$.  Isotropy is exactly
$\E R^2=d$.  Write
\[
 c_d=\frac{\Gamma(d/2)}{\sqrt\pi\,\Gamma((d-1)/2)},\qquad
 Q_R=\frac{\E R}{(d-1)\E R^{-1}}.
\]

\begin{theorem}[spherical saturation with exact asymptotic constants]
\label{thm:globalspherical}
Let $d\ge2$, $R>0$ almost surely, $\E R^2=d$ and $\E R^{-1}<\infty$.  Let $h\ge0$ have
$0<m_0,m_2<\infty$.  For $d=2$, additionally assume $h(t)\le Ce^{-a|t|}$; no additional condition is needed
when $d\ge3$.  Then
\begin{align}
 \alpha_{h,X}(r)&\sim \frac{c_d\E R}{d-1}\frac{m_0}{r},&
 \beta_{h,X}(r)&\sim c_d\E R^{-1}\frac{m_2}{r^3},\label{eq:sphericalasymp}\\
 \kappa_{h,X}(r)&\sim Q_R\frac{m_0}{m_2}r^2.\label{eq:sphericalkappa}
\end{align}
For $h_p=\sigma'^p$ and every $p>0$,
\begin{equation}
 \kappa_{p,X}(r)\sim\frac{Q_R}{2\psi_1(p)}r^2.
 \label{eq:sphericallogistic}
\end{equation}
\end{theorem}

\begin{proof}
Put $T=u^\top U$.  Its density is
$c_d(1-t^2)^{(d-3)/2}{\bf1}_{(-1,1)}(t)$, and conditional on $T=t$ every unit transverse coordinate has second
moment $(1-t^2)/(d-1)$.  Changing variables $s=rRT$ gives the exact identities
\begin{align}
 r\alpha_{h,X}(r)
 &=\frac{c_d}{d-1}\E\!\left[R\!\int_{|s|<rR}\!h(s)
 \left(1-\frac{s^2}{r^2R^2}\right)^{(d-1)/2}ds\right],\label{eq:sphericalexacta}\\
 r^3\beta_{h,X}(r)
 &=c_d\E\!\left[R^{-1}\!\int_{|s|<rR}\!s^2h(s)
 \left(1-\frac{s^2}{r^2R^2}\right)^{(d-3)/2}ds\right].\label{eq:sphericalexactb}
\end{align}
The first integrand is dominated by $Rm_0$.  For $d\ge3$ the second is dominated by $R^{-1}m_2$.
For $d=2$, splitting its inner integral at $|s|=rR/2$ gives a uniform bound from exponential localization;
the boundary contribution is bounded by a constant times $(rR)^3e^{-arR/2}$.  Dominated convergence proves
\eqref{eq:sphericalasymp}.  The ratio and the identity $m_0(p)/m_2(p)=1/[2\psi_1(p)]$ prove the rest.
\end{proof}

The factor $Q_R$ is the conditional transverse variance on the transition hyperplane.  It is distribution
sensitive and satisfies $0<Q_R\le d/(d-1)$.  Three closed examples are
\begin{center}
\begin{tabular}{lcc}
\toprule
isotropic spherical law & radius & $Q_R$\\
\midrule
$N(0,I_d)$ & $\chi_d$ & $1$\\
uniform on $\sqrt dS^{d-1}$ & $\sqrt d$ & $d/(d-1)$\\
isotropic Student $t_\nu$, $\nu>2$ & $\sqrt{\nu-2}\,\chi_d/\chi_\nu$ & $(\nu-2)/(\nu-1)$\\
\bottomrule
\end{tabular}
\end{center}
The last line follows from the adjacent chi-moment identities.  If $X=AS$ with $A$ invertible and $S$ isotropic
spherical, then exactly
\[
 H_{h,X}(w)=A H_{h,S}(\lVert A^\top w\rVert,A^\top w/\lVert A^\top w\rVert)A^\top.
\]
Thus the theorem transfers to elliptical inputs after whitening; in Euclidean coordinates it is a generalized
eigenvalue statement relative to $AA^\top$, not generally a literal radial--tangential eigendecomposition.

We now make the inverse-radius boundary exact.  The following elementary Abelian lemma is the only transition
calculus needed.  It is stated separately so that the critical logarithm and its coefficient can be checked without
appealing to an unnamed regular-variation theorem.

\begin{lemma}[spherical transition kernel]
\label{lem:transitionkernel}
Let $d\ge2$, let $h:\R\to[0,\infty)$ be bounded, and suppose
$0<m_2=\int_\R s^2h(s)\,ds<\infty$.  For $q>0$ define
\[
 J_{d,h}(q)=\int_{-1}^1t^2h(qt)(1-t^2)^{(d-3)/2}\,dt.
\]
When $d=2$, additionally assume $h(s)\le Ce^{-\lambda|s|}$; no extra condition is needed for $d\ge3$.  Then
\begin{equation}
 q^3J_{d,h}(q)\longrightarrow m_2,
 \qquad
 \int_0^xq^2J_{d,h}(q)\,dq\sim m_2\log x.
 \label{eq:kernellimits}
\end{equation}
For every $0<a<1$, writing $M_{a+1}(h)=\int_\R|s|^{a+1}h(s)\,ds$, one moreover has
\begin{equation}
 \int_0^\infty q^{a+1}J_{d,h}(q)\,dq
 =\frac12B\!\left(\frac{1-a}{2},\frac{d-1}{2}\right)M_{a+1}(h)<\infty.
 \label{eq:kernelfactor}
\end{equation}
\end{lemma}

\begin{proof}
For $d\ge3$, the change $s=qt$ gives
\[
 q^3J_{d,h}(q)=\int_{-q}^q s^2h(s)
 \left(1-\frac{s^2}{q^2}\right)^{(d-3)/2}ds.
\]
The integrand is bounded by $s^2h(s)$ and converges pointwise to it, proving the first limit by dominated
convergence.  For $d=2$, split the last integral at $|s|=q/2$.  Dominated convergence applies on the inner
part, while the outer part is $O(q^3e^{-\lambda q/2})$ after returning to the $t$ variable.  The same split gives
$J_{2,h}(q)=O(q^{-3})$.  Hence in every dimension $q^2J_{d,h}(q)=m_2/q+o(q^{-1})$, whose logarithmic integral proves the second limit.
For \eqref{eq:kernelfactor}, Tonelli's theorem and $s=q|t|$ give
\begin{align*}
 \int_0^\infty q^{a+1}J_{d,h}(q)\,dq
 &=M_{a+1}(h)\int_0^1t^{-a}(1-t^2)^{(d-3)/2}\,dt\\
 &=\frac12B\!\left(\frac{1-a}{2},\frac{d-1}{2}\right)M_{a+1}(h).
\end{align*}
Boundedness of $h$ controls the origin, while $M_{a+1}(h)<\infty$ follows from $a+1<2$ and the assumed
second moment.  This also proves finiteness directly.
\end{proof}

\begin{theorem}[exact three-phase inverse-radius boundary]
\label{thm:radialphase}
Let $d\ge2$, $X=RU$ be isotropic spherical, and let the positive radius have a density satisfying
\begin{equation}
 f_R(\rho)=c\rho^{a-1}(1+o(1))\quad(\rho\downarrow0),
 \qquad c>0,\ a>0,
 \label{eq:radialregularvariation}
\end{equation}
and $\E R^2=d$.  Let $h:\R\to[0,\infty)$ be bounded with $0<m_0,m_2<\infty$; when $d=2$, additionally
assume exponential localization as in Lemma~\ref{lem:transitionkernel}.  Put
\[
 A_{d,h,R}=\frac{c_d\E R}{d-1}m_0,
 \qquad
 D_{a,d,h}=\frac{c_dc}{2}B\!\left(\frac{1-a}{2},\frac{d-1}{2}\right)M_{a+1}(h)
 \quad(0<a<1).
\]
No assumption on the density away from zero is imposed beyond the displayed isotropic second moment.  Then
\begin{equation}
 \alpha_{h,X}(r)\sim A_{d,h,R}r^{-1},
 \label{eq:phasealpha}
\end{equation}
and the radial eigenvalue has the exact three-phase law
\begin{equation}
 \beta_{h,X}(r)\sim
 \begin{cases}
 D_{a,d,h}\,r^{-(a+2)},&0<a<1,\\[2mm]
 c_dc\,m_2\,r^{-3}\log r,&a=1,\\[2mm]
 c_d\E R^{-1}m_2\,r^{-3},&a>1.
 \end{cases}
 \label{eq:exactthreephase}
\end{equation}
Consequently
\begin{equation}
 \kappa_{h,X}(r)\sim
 \begin{cases}
 (A_{d,h,R}/D_{a,d,h})r^{a+1},&0<a<1,\\[1mm]
 \dfrac{\E R\,m_0}{(d-1)c m_2}\dfrac{r^2}{\log r},&a=1,\\[3mm]
 \dfrac{\E R}{(d-1)\E R^{-1}}\dfrac{m_0}{m_2}r^2,&a>1.
 \end{cases}
 \label{eq:exactthreephasekappa}
\end{equation}
Thus $\E R^{-1}<\infty$ is exactly the quadratic-phase boundary inside the class
\eqref{eq:radialregularvariation}.
\end{theorem}

\begin{proof}
The identity \eqref{eq:sphericalexacta} is dominated by $Rm_0$.  Since $\E R<\infty$ follows from isotropy,
dominated convergence proves \eqref{eq:phasealpha} for every $a>0$.

For the radial direction, direct conditioning on $R=\rho$ and then putting $q=r\rho$ gives
\begin{equation}
 \beta_{h,X}(r)=\frac{c_d}{r^3}\int_0^\infty f_R(q/r)q^2J_{d,h}(q)\,dq.
 \label{eq:phaseexactq}
\end{equation}
Fix $\eta>0$ small enough that \eqref{eq:radialregularvariation} supplies
$f_R(\rho)\le C\rho^{a-1}$ on $(0,\eta)$.  The contribution of $R\ge\eta$ obeys, directly from
\eqref{eq:sphericalexactb},
\begin{equation}
 r^3\beta_{h,X}^{(R\ge\eta)}(r)\le C_{d,h}/\eta.
 \label{eq:phasetail}
\end{equation}
For $d\ge3$ one may take $C_{d,h}=c_dm_2$.  For $d=2$, the split in
Lemma~\ref{lem:transitionkernel} proves that $q^3J_{2,h}(q)$ is uniformly bounded; hence one may take
\[
 C_{2,h}=c_d\sup_{q>0}q^3J_{2,h}(q)
\]
directly.

Suppose first that $0<a<1$.  On $q<\eta r$, the integrand of $r^{a+2}\beta/c_d$ converges pointwise to
$cq^{a+1}J_{d,h}(q)$ and is bounded by $Cq^{a+1}J_{d,h}(q)$.  Lemma~\ref{lem:transitionkernel} makes this
majorant integrable.  Dominated convergence and \eqref{eq:phasetail}, whose scaled contribution is
$O(r^{a-1})$, yield
\[
 r^{a+2}\beta_{h,X}(r)\longrightarrow
 c_dc\int_0^\infty q^{a+1}J_{d,h}(q)\,dq=D_{a,d,h}.
\]

At $a=1$, for every $\epsilon>0$ one may decrease $\eta$ so that
$c-\epsilon\le f_R(\rho)\le c+\epsilon$ on $(0,\eta)$.  Therefore the small-radius part in
\eqref{eq:phaseexactq} is squeezed between $c\pm\epsilon$ times
$\int_0^{\eta r}q^2J_{d,h}(q)\,dq$.  The second limit in \eqref{eq:kernellimits} and the $O(1)$ tail
bound \eqref{eq:phasetail}, followed by $\epsilon\downarrow0$, prove
$r^3\beta/(\log r)\to c_dcm_2$.

Finally, $a>1$ makes $\E R^{-1}<\infty$: local integrability follows from
$\rho^{-1}f_R(\rho)\asymp\rho^{a-2}$, and the part above any fixed radius is automatic.  Theorem
\ref{thm:globalspherical} proves the last line.  Dividing \eqref{eq:phasealpha} by
\eqref{eq:exactthreephase} proves \eqref{eq:exactthreephasekappa}.
\end{proof}

For a concrete isotropic family in every phase, let
\[
 Q\sim\operatorname{Gamma}(a,1),\qquad R=\sqrt{d/[a(a+1)]}\,Q.
\]
Then $\E R^2=d$ and
$c=\{\Gamma(a)[d/(a(a+1))]^{a/2}\}^{-1}$.  Figure~\ref{fig:radialphases} evaluates the defining integrals
for $a\in\{1/2,1,2\}$ and normalizes by the constants in \eqref{eq:exactthreephase}; it is a diagnostic of all
three limits, not evidence used by the proof.

\begin{figure}[H]
 \centering
 \includegraphics[width=.88\linewidth]{../figures/radial_phase_transition.pdf}
 \caption{Exact spherical phase normalizations for $d=3$, $h=\sigma'$, and isotropic Gamma radii.  Every curve
 tends to one under its theorem-specific normalization; the slow approach, especially in the critical $a=1$
 regime, is part of the predicted asymptotics.  Values are deterministic adaptive quadrature.}
 \label{fig:radialphases}
\end{figure}

Covariance isotropy alone still does not force any phase law.  For independent Rademacher coordinates,
$u=e_1$ and even $h$, one has $H_h(r,u)=h(r)I_d$ and hence $\kappa=1$: here $X_1^2=1$, all mixed expectations
vanish by sign symmetry, and $h(rX_1)=h(r)$.

The radius--direction representation is classical \cite{cambanis1981}; non-Gaussian reference exponents also
appear in regular-design logistic theory \cite{chardon2024}.  The contribution here is the exact $Q_R$ constant,
its all-$p$ logistic specialization, and the sharp inverse-radius boundary within the stated regularly varying
class, not the spherical representation or the
first occurrence of non-Gaussian $r^{-1}$/$r^{-3}$ scales.

\section{Spherical and finite-scale Schwarzian bridges}

The matrix-order and saturation calculations meet exactly at an inflection point.  Gaussian independence makes the
full saturation curve one-dimensional, but it is not responsible for the local bridge: rotational symmetry and
suitable moments suffice.  The fourth moment identifies the coefficient below; the stated sixth moment controls
its $O(r^4)$ remainder.

\begin{theorem}[spherical local order--anisotropy bridge]
\label{thm:bridge}
Let $d\ge2$ and let $X\in\R^d$ be isotropic and spherically symmetric with
$\E\lVert X\rVert^6<\infty$.  Let $g\in C^5(\R)$ satisfy $g'>0$ and $g''(0)=0$.
For $p>0$, suppose $h_p=(g')^p$ has bounded fourth derivative.  For $u\in S^{d-1}$ put
$Z=u^\top X$ and
\[
 H_{p,X}(r,u)=\E[h_p(rZ)XX^\top]
 =\alpha_{p,X}(r)(I-uu^\top)+\beta_{p,X}(r)uu^\top ,
 \qquad \kappa_{p,X}=\frac{\alpha_{p,X}}{\beta_{p,X}}.
\]
Then, with
\[
 q_X:=\frac{\E\lVert X\rVert^4}{d(d+2)},
\]
one has
\begin{equation}
 \lim_{r\downarrow0}\frac{\kappa_{p,X}(r)-1}{r^2}
 =-q_Xp\,Sg(0)
 =-\frac{6q_Xp}{g'(0)^2}
   \lim_{\delta\to0}\frac{\det L_g(0,\delta)}{\delta^2}.
 \label{eq:bridge}
\end{equation}
In particular, $Sg(0)<0$ makes $L_g(0,\delta)$ indefinite and
$\kappa_{p,X}(r)>1$ at all sufficiently small nonzero scales.  Standard Gaussian input has $q_X=1$.
\end{theorem}

\begin{proof}
Rotations fixing $u$ force the displayed two-eigenspace form, with
\[
 \beta_{p,X}(r)=\E[Z^2h_p(rZ)],\qquad
 \alpha_{p,X}(r)=\frac1{d-1}
 \E[(\lVert X\rVert^2-Z^2)h_p(rZ)].
\]
Choose a unit $v\perp u$.  Spherical symmetry and isotropy give
\[
 \E Z^4=3q_X,\qquad \E[Z^2(v^\top X)^2]=q_X.
\]
Taylor's formula, cancellation of odd terms under $X\stackrel d=-X$, and the sixth-moment assumption therefore give
\[
 \alpha_{p,X}(r)=h_p(0)+\frac{q_Xh_p''(0)}2r^2+O(r^4),\qquad
 \beta_{p,X}(r)=h_p(0)+\frac{3q_Xh_p''(0)}2r^2+O(r^4).
\]
Taking the ratio yields
$\kappa_{p,X}(r)=1-q_Xh_p''(0)r^2/h_p(0)+O(r^4)$.
Since $\log h_p=p\log g'$ and $g''(0)=0$,
\[
 \frac{h_p''(0)}{h_p(0)}=p\frac{g'''(0)}{g'(0)}=pSg(0).
\]
The adjacent-point expansion \eqref{eq:localschwarz} gives the remaining equality.
\end{proof}

The equality of limits has a computable finite-scale version.  The constants below are deliberately explicit rather
than optimized.

\begin{theorem}[explicit finite-scale bridge]
\label{thm:finitebridge}
Under Gaussian input, retain the hypotheses of Theorem~\ref{thm:bridge} and write
\[
 a=h_p(0),\quad b=h_p''(0),\quad M_4=\lVert h_p^{(4)}\rVert_\infty.
\]
Choose $r_0>0$ so that
\begin{equation}
 \frac{3|b|}{2}r_0^2+\frac{5M_4}{8}r_0^4\le\frac a2,
 \label{eq:r0condition}
\end{equation}
and set
\[
 C_h=\frac{3M_4}{2a}+3\left(\frac ba\right)^2
       +\frac{5|b|M_4}{4a^2}r_0^2.
\]
For a chosen $\delta_0>0$, put
\[
 A=g'(0),\ c=g'''(0),\ d_4=g^{(4)}(0),\
 M_5=\sup_{|x|\le\delta_0}|g^{(5)}(x)|,
\]
\[
 K=\frac{|c|}{6}+\frac{|d_4|\delta_0}{24}
      +\frac{M_5\delta_0^2}{120},\qquad
 D_g=\frac{|Ad_4|}{12}+\frac{7AM_5\delta_0}{120}+\delta_0K^2.
\]
Then for $0<r\le r_0$ and $0<|\delta|\le\delta_0$,
\begin{align}
 \left|\kappa_p(r)-1+pSg(0)r^2\right|&\le C_hr^4,
 \label{eq:quantkappa}\\
 \left|\det L_g(0,\delta)-\frac{A^2Sg(0)}6\delta^2\right|
 &\le D_g|\delta|^3,
 \label{eq:quantloewner}\\
 \left|\frac{\kappa_p(r)-1}{r^2}
 +\frac{6p}{A^2}\frac{\det L_g(0,\delta)}{\delta^2}\right|
 &\le C_hr^2+\frac{6pD_g}{A^2}|\delta|.
 \label{eq:directfinitebridge}
\end{align}
Thus the order defect predicts anisotropy at finite, independently chosen scales with a certified error bar.
\end{theorem}

\begin{proof}
Gaussian Taylor remainders give
\[
 \alpha_p=a+\frac b2r^2+R_\alpha,\quad |R_\alpha|\le\frac{M_4}{8}r^4,
 \qquad
 \beta_p=a+\frac{3b}{2}r^2+R_\beta,\quad |R_\beta|\le\frac{5M_4}{8}r^4.
\]
Condition \eqref{eq:r0condition} makes $\beta_p\ge a/2$.  Subtracting
$1-(b/a)r^2$ before dividing by $\beta_p$ yields
\[
 \left|\kappa_p-1+\frac ba r^2\right|
 \le\left[
 \frac{3M_4}{2a}+3\left(\frac ba\right)^2
 +\frac{5|b|M_4}{4a^2}r_0^2\right]r^4,
\]
which is \eqref{eq:quantkappa} because $b/a=pSg(0)$.

Taylor expansion at zero gives
\[
 g'(\delta)=A+\frac c2\delta^2+\frac{d_4}{6}\delta^3+R_1,\quad
 \frac{g(\delta)-g(0)}\delta
 =A+\frac c6\delta^2+\frac{d_4}{24}\delta^3+R_2,
\]
where $|R_1|\le M_5|\delta|^4/24$ and
$|R_2|\le M_5|\delta|^4/120$.  Writing the determinant as
$A(g'(\delta)-A)-2A(V-A)-(V-A)^2$, with
$V=(g(\delta)-g(0))/\delta$, gives \eqref{eq:quantloewner} with the displayed $D_g$.
Finally, divide \eqref{eq:quantkappa} by $r^2$, divide
\eqref{eq:quantloewner} by $\delta^2$, multiply the latter by $6p/A^2$, and use the triangle inequality.
\end{proof}

For the standard logistic sigmoid, $S\sigma=-1/2$.  Theorem~\ref{thm:bridge} therefore predicts the exact initial
coefficient $p/2$ for every $h_p=\sigma'^p$, recovered globally below.  Equation
\eqref{eq:directfinitebridge} strengthens the conceptual coefficient match into a falsifiable finite-scale
comparison.

For a general $h$, $\kappa_h=\alpha_h/\beta_h$ in \eqref{eq:univasymp} is an anisotropy ratio; it is the spectral
condition number only after the eigenvalue ordering is known.  The two powers are geometric.  At large $r$, only a
score slab of width $O(r^{-1})$ remains unsaturated, producing
$\alpha\asymp r^{-1}$.  A radial perturbation carries an additional factor $Z^2=O(r^{-2})$ inside that slab,
producing $\beta\asymp r^{-3}$.  Both modes are supported by the near-boundary slab; the radial mode is weaker
because it carries the additional $Z^2$ factor.

\section{All-power logistic saturation laws}

For every real $p>0$ set $h_p(t)=\sigma'(t)^p$.  The cases $p=1$ and $p=2$ are Bernoulli Fisher and
squared-output curvature, respectively, but the analytic law is not restricted to those observation models.  The
Gamma transform below is classical generalized-logistic distribution theory \cite[eqs.~(2.3)--(2.5)]{ojo2004};
our use of it is geometric.

\begin{lemma}[closed sensitivity moments for every $p>0$]
\label{lem:moments}
Let $\psi_j$ denote the polygamma function of order $j$.  Then
\begin{equation}
 \widehat h_p(k):=\int_\R e^{ikt}h_p(t)\,dt
 =\frac{\Gamma(p+ik)\Gamma(p-ik)}{\Gamma(2p)},
 \label{eq:fourier}
\end{equation}
and
\begin{equation}
 m_0(p)=\frac{\Gamma(p)^2}{\Gamma(2p)},\qquad
 m_2(p)=2\psi_1(p)m_0(p),\qquad
 m_4(p)=\bigl(12\psi_1(p)^2+2\psi_3(p)\bigr)m_0(p).
 \label{eq:allpmoments}
\end{equation}
In particular,
\begin{center}
\begin{tabular}{c@{\qquad}ccc}
\toprule
profile & $m_0$ & $m_2$ & $m_4$\\
\midrule
$h_1=\sigma'$ & $1$ & $\pi^2/3$ & $7\pi^4/15$\\[2pt]
$h_2=\sigma'^2$ & $1/6$ & $(\pi^2-6)/18$ & $7\pi^4/90-2\pi^2/3$\\
\bottomrule
\end{tabular}
\end{center}
\end{lemma}

\begin{proof}
With $x=e^t$,
\[
 h_p(t)=\frac{x^p}{(1+x)^{2p}},\qquad
 \widehat h_p(k)=\int_0^\infty\frac{x^{p+ik-1}}{(1+x)^{2p}}\,dx
 =B(p+ik,p-ik),
\]
which proves \eqref{eq:fourier}.  Since $h_p(t)=O(e^{-p|t|})$, all polynomial moments are finite and
differentiation under the Fourier integral is justified.  Differentiating its logarithm at zero gives
$(\log\widehat h_p)''(0)=-2\psi_1(p)$ and
$(\log\widehat h_p)^{(4)}(0)=2\psi_3(p)$.  Since
$m_2=-\widehat h_p''(0)$ and $m_4=\widehat h_p^{(4)}(0)$,
\eqref{eq:allpmoments} follows.  The table uses the standard integer polygamma values.
\end{proof}

\begin{theorem}[strict all-power logistic anisotropy]
\label{thm:monotone}
For every $p>0$, $\kappa_p(0)=1$, $\kappa_p(r)>1$ for $r>0$, and $\kappa_p$ is strictly increasing on
$(0,\infty)$.  At the two endpoints,
\begin{equation}
 \kappa_p(r)=1+\frac p2r^2-\frac p8r^4+\frac{p(p+1)}{16}r^6+O(r^8)
 \quad(r\downarrow0),
 \label{eq:smallallp}
\end{equation}
and
\begin{equation}
 \kappa_p(r)=\frac{r^2}{2\psi_1(p)}
 +1+\frac{\psi_3(p)}{4\psi_1(p)^2}+O(r^{-2})
 \quad(r\to\infty).
 \label{eq:logisticconstants}
\end{equation}
Thus $\kappa_1(r)\sim3r^2/\pi^2$ and
$\kappa_2(r)\sim3r^2/(\pi^2-6)$.
\end{theorem}

\begin{proof}
Normalize the density $q_{p,r}(z)\propto h_p(rz)\varphi(z)$.  Then
$\beta_p(r)/\alpha_p(r)=\E_{q_{p,r}}Z^2$.  If $r_2>r_1$, the derivative on $z>0$ of the log likelihood ratio is
\[
 \frac d{dz}\log\frac{h_p(r_2z)}{h_p(r_1z)}
 =-p\bigl[r_2\tanh(r_2z/2)-r_1\tanh(r_1z/2)\bigr]<0.
\]
Thus the distribution of $|Z|$ under $q_{p,r}$ decreases strictly in monotone-likelihood-ratio order as $r$
increases, and its strictly increasing statistic $Z^2$ has decreasing expectation.  Hence $\alpha_p/\beta_p$
increases strictly.  At $r=0$ isotropy gives one.  Lemma~\ref{lem:moments} inserted into
Theorem~\ref{thm:universal}, with one further term in
$e^{-t^2/(2r^2)}$, gives the large-$r$ statement.  More explicitly, $m_6(p)<\infty$ by the exponential tail,
and Taylor remainder bounds for $e^{-x}$ give
\[
 \alpha_p(r)=\frac{\varphi(0)}r\left(m_0-\frac{m_2}{2r^2}
 +\frac{m_4}{8r^4}+O(r^{-6})\right),\qquad
 \beta_p(r)=\frac{\varphi(0)}{r^3}\left(m_2-\frac{m_4}{2r^2}
 +O(r^{-4})\right).
\]
Their ratio yields \eqref{eq:logisticconstants}.  Finally,
\[
 h_p(t)=4^{-p}\left[1-\frac p4t^2+\frac{p(1+3p)}{96}t^4
 -\frac{p(15p^2+15p+4)}{5760}t^6+O(t^8)\right].
\]
For each fixed $p>0$, every derivative of $h_p$ is bounded.  Taylor's theorem therefore bounds the displayed
remainder globally by $\lVert h_p^{(8)}\rVert_\infty |t|^8/8!$; after setting $t=rZ$, this is integrable in both
$\alpha_p$ and $\beta_p$ because the Gaussian has a finite tenth moment.  Inserting the Gaussian moments through
$\E Z^8=105$ in the coefficients and taking the ratio gives \eqref{eq:smallallp}.
\end{proof}

Figure~\ref{fig:saturation} shows the two statistically canonical members of the all-power family.  The solid and
dashed curves are direct Gaussian integrals; the dotted curves are the asymptotic constants in
\eqref{eq:logisticconstants}.

\begin{figure}[H]
 \centering
 \includegraphics[width=\textwidth]{../figures/saturation_law.pdf}
 \caption{Exact curvature/information eigenvalues and their ratio.  Saturation leaves $d-1$ tangential directions
 at order $r^{-1}$ but pushes the radial direction to order $r^{-3}$.  The quadratic condition-number law is not a
 fitted exponent; it follows from Theorem~\ref{thm:universal}.}
 \label{fig:saturation}
\end{figure}

The cancellation of every $p^2$ contribution in the fourth-order ratio coefficient is worth noting: the full
all-power family has the correction $-pr^4/8$, followed by the first nonlinear-in-$p$ term
$p(p+1)r^6/16$.

\section{Finite-sample resolution of the weighted Gram matrix}
\label{sec:finite}

The population eigenvalues alone do not say when a sample resolves the radial mode.  We now answer that question
for the fixed-oracle matrix
\[
 \widehat H_{h,n}(r,u)=\frac1n\sum_{i=1}^n h(ru^\top X_i)X_iX_i^\top,
 \qquad X_i\stackrel{\rm iid}{\sim}N(0,I_d).
\]
Prior self-concordant analyses already give bilateral constant-factor empirical-Hessian sandwiches, including a
uniform result for Gaussian logistic regression \cite[Eq.~(92)]{ostrovskii2021}; fixed-point standardized-Hessian
concentration is also available under generic matrix-Bernstein conditions \cite{fisher2023}.  Chardon, Lerasle and
Mourtada prove a uniform one-sided lower bound for the empirical Bernoulli Hessian ($p=1$) under the sufficient
condition $n\gtrsim r(d+t)$ \cite[Theorem 6]{chardon2024}.  We do not claim the first bilateral concentration or
the first finite-sample lower bound.  The result below instead gives tunable two-sided relative error for a fixed
teacher, applies to a general bounded profile and hence every $p>0$, separates radial, tangential and cross errors,
and returns an eigenspace angle.

Put $m=d-1$, $Z=u^\top X$, $Y=X-Zu$, and $W=h(rZ)$.  Define
\[
 \alpha=\E W,\quad \beta=\E[WZ^2],\quad
 \gamma_j=\E[W^2Z^{2j}]\quad(j=0,1,2),
\]
and the envelopes
\[
 H=\lVert h\rVert_\infty,\qquad
 K_1=\sup_z z^2h(rz)^2,\qquad K_2=\sup_z z^2h(rz).
\]

\begin{theorem}[explicit relative empirical spectrum]
\label{thm:empirical}
Assume $0\le h\le H$, $K_1,K_2<\infty$, and $\alpha,\beta>0$.  Let
$0<\delta<1$ and set
\[
 t=\log(12/\delta),\qquad u_0=m\log9+t,
\]
\[
 S_0^+=n\gamma_0+\sqrt{2nH^2\gamma_0t}+\frac{H^2t}{3},\qquad
 S_1^+=n\gamma_1+\sqrt{2nK_1\gamma_1t}+\frac{K_1t}{3},
\]
\begin{align}
 e_T&=\sqrt{\frac{2\gamma_0t}{n}}+\frac{Ht}{3n}
 +\frac4n\bigl(\sqrt{S_0^+u_0}+Hu_0\bigr),\label{eq:et}\\
 e_R&=\sqrt{\frac{2\gamma_2t}{n}}+\frac{K_2t}{3n},\qquad
 q_n=\frac{\sqrt{S_1^+}}n\bigl(\sqrt m+\sqrt{2t}\bigr).\label{eq:erq}
\end{align}
Then, with probability at least $1-\delta$, in the radial--tangential basis
\[
 \widehat H_{h,n}=\begin{pmatrix}C&b_n\\b_n^\top&c_n\end{pmatrix}
\]
satisfies simultaneously
\begin{equation}
 \lVert C-\alpha I_m\rVert_{\rm op}\le e_T,\qquad
 |c_n-\beta|\le e_R,\qquad \lVert b_n\rVert_2\le q_n.
 \label{eq:blockevent}
\end{equation}
Consequently, for
\begin{equation}
 \varepsilon_n:=\max\left\{\frac{e_T}{\alpha},\frac{e_R}{\beta}\right\}
 +\frac{q_n}{\sqrt{\alpha\beta}},
 \label{eq:epsn}
\end{equation}
one has the two-sided relative L\"owner bound
\begin{equation}
 (1-\varepsilon_n)H_h(r,u)\preceq\widehat H_{h,n}(r,u)
 \preceq(1+\varepsilon_n)H_h(r,u).
 \label{eq:relativegram}
\end{equation}
If $G:=\alpha-\beta-e_T-e_R>0$, the empirical matrix has a unique eigenvalue
$\widehat\beta$ below the tangential block spectrum, and its unit eigenvector $\widehat u$ obeys
\begin{equation}
 \tan\angle(\widehat u,u)\le\frac{q_n}{G},\qquad
 |\widehat\beta-\beta|\le e_R+\frac{q_n^2}{G}.
 \label{eq:angle}
\end{equation}
Every other ordered eigenvalue differs from $\alpha$ by at most
$\max\{e_T,e_R\}+q_n$.
\end{theorem}

\begin{proof}
Write
\[
 C=\frac1n\sum_iW_iY_iY_i^\top,\qquad
 b_n=\frac1n\sum_iW_iZ_iY_i,\qquad
 c_n=\frac1n\sum_iW_iZ_i^2.
\]
Scalar Bernstein controls $\overline W-\alpha$ and, one-sidedly,
$\sum_iW_i^2\le S_0^+$.  Conditional on the weights, for a fixed unit $v$,
\[
 v^\top\sum_iW_i(Y_iY_i^\top-I_m)v=\sum_iW_i(G_i^2-1),
 \qquad G_i\stackrel{\rm iid}{\sim}N(0,1).
\]
The exact Gaussian moment-generating function gives
\[
 \left|\sum_iW_i(G_i^2-1)\right|
 \le2\sqrt{u_0\sum_iW_i^2}+2Hu_0
\]
outside conditional probability $2e^{-u_0}$.  A $1/4$-net of $S^{m-1}$ has at most $9^m$ points and
$\lVert A\rVert_{\rm op}\le2\max_{v\ {\rm in\ net}}|v^\top Av|$, proving the first bound in
\eqref{eq:blockevent}.  Bernstein applied to $W_iZ_i^2\in[0,K_2]$ proves the radial bound.

Conditionally on the $Z_i$,
\[
 b_n\sim N\!\left(0,\frac{\sum_iW_i^2Z_i^2}{n^2}I_m\right).
\]
A second one-sided Bernstein bound gives $\sum_iW_i^2Z_i^2\le S_1^+$, and the Gaussian norm tail proves the
cross bound.  The listed failures total less than $9e^{-t}<\delta$.

Conjugating $\widehat H_{h,n}-H_h$ by
$H_h^{-1/2}=\operatorname{diag}(\alpha^{-1/2}I_m,\beta^{-1/2})$ shows that its norm is at most
\eqref{eq:epsn}; this is equivalent to \eqref{eq:relativegram}.  If $G>0$, then
$c_n<\lambda_{\min}(C)$.  Interlacing leaves one eigenvalue below $C$.
The eigenvector equation and the Schur complement give respectively
$\tan\angle(\widehat u,u)\le q_n/G$ and
$0\le c_n-\widehat\beta\le q_n^2/G$.  Weyl's inequality gives the remaining assertion.
\end{proof}

For $h_p=\sigma'^p$, all quantities in Theorem~\ref{thm:empirical} are one-dimensional and computable.
The elementary bound $\sigma'(s)\le e^{-|s|}$ gives, for every $p>0$,
\begin{equation}
 H=4^{-p},\quad K_1\le\frac{e^{-2}}{p^2r^2},\quad
 K_2\le\frac{4e^{-2}}{p^2r^2},
 \label{eq:envelopes}
\end{equation}
and, with $\varphi(0)=(2\pi)^{-1/2}$,
\begin{equation}
 \gamma_0\le\frac{\varphi(0)}{pr},\qquad
 \gamma_1\le\frac{\varphi(0)}{2p^3r^3},\qquad
 \gamma_2\le\frac{3\varphi(0)}{2p^5r^5}.
 \label{eq:gammabounds}
\end{equation}
Indeed, after $s=rZ$, use $h_p(s)^2\le e^{-2p|s|}$ and
$\int_\R |s|^je^{-2p|s|}ds=2j!/(2p)^{j+1}$ for $j=0,2,4$.

\begin{corollary}[effective sample size for every logistic power]
\label{cor:effective}
Let
\[
 R_p^2=\max\left\{\frac{3m_2(p)}{m_0(p)},\frac{m_4(p)}{m_2(p)}\right\}.
\]
For $r\ge R_p$, $0<\varepsilon\le1/2$ and $0<\delta<1$, the explicit condition
$\varepsilon_n\le\varepsilon$ in \eqref{eq:epsn}, with
\eqref{eq:envelopes}--\eqref{eq:gammabounds} substituted, guarantees
\eqref{eq:relativegram}.  In particular, for each fixed $p>0$ there are finite computable constants $C_p,C'_p$
such that
\begin{equation}
 n\ge C_p\,\frac{r\,[d+\log(12/\delta)]}{\varepsilon^2}
 \label{eq:samplesufficient}
\end{equation}
is sufficient and, after increasing $C_p$ if necessary,
$e_T+e_R\le(\alpha-\beta)/2$.  Hence $G\ge(\alpha-\beta)/2$ and, on the same event,
\[
 \sin\angle(\widehat u,u)
 \le C'_p\!\left(\sqrt{\frac{d+\log(1/\delta)}{nr}}
 +\frac{\sqrt{[d+\log(1/\delta)]\log(1/\delta)}}{n}\right).
\]
Thus the effective number of transition-layer observations is $n/r$: relative resolution of a population
eigenvalue of order $r^{-3}$ does not itself require $r^3$ samples.
\end{corollary}

\begin{proof}
The brackets in Theorem~\ref{thm:universal} and the definition of $R_p$ give
\[
 \alpha\ge\frac{\varphi(0)m_0(p)}{2r},\qquad
 \beta\ge\frac{\varphi(0)m_2(p)}{2r^3},\qquad
 \alpha-\beta\ge\frac{\varphi(0)m_0(p)}{2r}.
\]
Insert these and \eqref{eq:envelopes}--\eqref{eq:gammabounds} in
\eqref{eq:et}--\eqref{eq:epsn}.  Each relative block term is bounded by a $p$-dependent constant times
$\sqrt{r[d+\log(12/\delta)]/n}+r[d+\log(12/\delta)]/n$.
The cross term has the same first order.  Increasing a finite constant $C_p$ makes their sum at most
$\varepsilon$ and also makes $e_T+e_R\le(\alpha-\beta)/2$.  Thus
$G\ge(\alpha-\beta)/2\ge\varphi(0)m_0(p)/(4r)$, and substitution of $q_n$ into
\eqref{eq:angle} gives the displayed bound with a finite computable $C'_p$.
\end{proof}

For the two statistical powers, the conservative majorant used in this proof gives fully numerical choices:
\begin{center}
\begin{tabular}{cccc}
\toprule
$p$ & $R_p$ & $C_p$ in \eqref{eq:samplesufficient} & $C'_p$ in the angle bound\\
\midrule
$1$ & $3.717183$ & $22929$ & $22.479249$\\
$2$ & $2.153017$ & $294162$ & $52.368963$\\
\bottomrule
\end{tabular}
\end{center}
These are sufficient, not optimized, constants.  The explicit-constants replay prints every intermediate
coefficient and checks that the resulting quadratic majorant is at most $\varepsilon$ for
$0<\varepsilon\le1/2$.

We next prove that the product scale in \eqref{eq:samplesufficient} is intrinsic for the stated bilateral loss.
Choose any finite $a_p>0$ satisfying
\[
 \int_{|s|>a_p}s^2h_p(s)\,ds\le\frac{m_2(p)}{16}.
\]
For
$r\ge\max\{\sqrt{m_4(p)/m_2(p)},\,4\varphi(0)a_p\}$, the empty-transition-slab event and Markov's inequality give
\begin{equation}
 \Pr\!\left\{\lambda_{\min}(\widehat H_{p,n})<\frac{\beta_p(r)}2\right\}
 \ge\frac12\exp\!\left(-\frac{4\varphi(0)a_pn}{r}\right).
 \label{eq:lower}
\end{equation}
Hence failure probability at most $\delta<1/2$ requires
\begin{equation}
 n\ge\frac{r}{4\varphi(0)a_p}\log\frac1{2\delta}.
 \label{eq:lowern}
\end{equation}
To verify \eqref{eq:lower}, condition on no $|Z_i|\le a_p/r$.  This event has probability at least the exponential
shown; the conditional mean radial entry is at most $\beta_p/4$, so with conditional probability at least $1/2$
that entry, and therefore the smallest eigenvalue, is below $\beta_p/2$.

The high-confidence half of the matching lower bound is an exact consequence of the next lemma.  Stating it
autonomously exposes every constant and removes any dependence on an implicit moderate-deviation argument.

\begin{lemma}[explicit Gaussian-product lower tail]
\label{lem:gaussianproducttail}
Let $A_i\ge0$ be iid with $0<\rho=\E A_i^2<\infty$ and
$\chi=\E A_i^4/\rho^2<\infty$.  Let $Y_i,Y_i'$ be independent standard Gaussians, independent of the $A_i$, and
put $D_n=n^{-1}\sum_{i=1}^nA_iY_iY_i'$.  For every $\varepsilon>0$,
\begin{equation}
 \Pr\{|D_n|>\varepsilon\}
 \ge \frac{3}{40(1+3\chi/n)}
 \exp\!\left(-\frac{2n\varepsilon^2}{\rho}\right).
 \label{eq:producttail}
\end{equation}
In particular, if $n\ge3\chi$, the leading constant is at least $3/80$.
\end{lemma}

\begin{proof}
Set $T=n^{-1}\sum_iA_i^2Y_i^2$.  Direct expansion gives
\[
 \E T=\rho,
 \qquad
 \E T^2=\rho^2\left(1-\frac1n+\frac{3\chi}{n}\right)
 \le\rho^2(1+3\chi/n).
\]
Paley--Zygmund therefore gives
$\Pr\{T\ge\rho/2\}\ge[4(1+3\chi/n)]^{-1}$.  Conditional on $(A_i,Y_i)_{i\le n}$,
$D_n\sim N(0,T/n)$.  The elementary bound
\begin{equation}
 2\Phi(-x)\ge\frac3{10}e^{-x^2},\qquad x\ge0,
 \label{eq:elementarygaussianlower}
\end{equation}
then proves \eqref{eq:producttail}.  For completeness, on $0\le x\le1$ use
$2\Phi(-x)\ge2\Phi(-1)>3/10$; on $x\ge1$ use the Mills bound
$\Phi(-x)\ge\varphi(x)x/(1+x^2)$, whose ratio to $e^{-x^2}$ is increasing from a value larger than $3/10$.
\end{proof}

\begin{theorem}[matching lower complexity for bilateral relative Loewner loss]
\label{thm:matchinglower}
Fix $p>0$.  There are finite constants $\widetilde R_p,c_p,\delta_p>0$ such that, for $r\ge\widetilde R_p$,
$d\ge3$, $0<\varepsilon\le1/2$ and $0<\delta\le\delta_p$, the implication
\begin{equation}
 \Pr\{(1-\varepsilon)H_p\preceq\widehat H_{p,n}\preceq(1+\varepsilon)H_p\}\ge1-\delta
 \quad\Longrightarrow\quad
 n\ge c_p\frac{r[d+\log(1/\delta)]}{\varepsilon^2}
 \label{eq:matchinglower}
\end{equation}
holds.  At constant confidence the dimension term has the explicit sufficient obstruction
\begin{equation}
 n<\frac{m_0(2p)}{16\varphi(0)m_0(p)^2}\,
 \frac{r(d-2)}{\varepsilon^2}
 \quad\Longrightarrow\quad
 \Pr\{\text{relative Loewner failure}\}\ge\frac1{1152},
 \label{eq:explicitlower}
\end{equation}
 provided $n\gamma_0\ge H^2$ and the finite-radius moment brackets hold.  The displayed coefficient equals
 $1/[96\varphi(0)]$ for $p=1$ and $9/[560\varphi(0)]$ for $p=2$.

The confidence constants can be made fully explicit.  Define
\begin{equation}
 L_p=\frac{m_0(2p)}{2\varphi(0)m_0(p)^2},
 \qquad
 K_p=\frac{4m_0(4p)}{\varphi(0)m_0(2p)^2},
 \qquad
 \delta_p^{\rm hc}=\frac12e^{-12\varphi(0)a_pK_p}.
 \label{eq:explicitconfidenceconstants}
\end{equation}
After enlarging $\widetilde R_p$ to ensure $r^2\ge m_2(2p)/m_0(2p)$, success necessarily implies the following
bound whenever $0<\delta<\min\{3/80,\delta_p^{\rm hc}\}$:
\begin{equation}
 n\ge \frac{L_p}{2}\frac r{\varepsilon^2}
 \log\frac{3}{80\delta}.
 \label{eq:explicithighconfidence}
\end{equation}
\end{theorem}

\begin{proof}
Restrict the relative inequality to the $m=d-1$ dimensional tangential block
$C=n^{-1}\sum_iW_iY_iY_i^\top$.  Put $\gamma_0=\E W^2$ and $S_0=\sum_iW_i^2$.  Since $W\le H$,
Paley--Zygmund gives
\[
 \Pr\{S_0\ge n\gamma_0/2\}\ge1/8\qquad(n\gamma_0\ge H^2).
\]
Conditional on the weights, $T_0=\sum_iW_i^2Y_{i1}^2$ has second moment at most $3S_0^2$, so
$\Pr\{T_0\ge S_0/2\mid W\}\ge1/12$.  Conditional on $(W,Y_1)$, the remaining $d-2$ entries of the first
off-diagonal column of $C$ are Gaussian with common variance $T_0/n^2$.  A final Paley--Zygmund bound gives
$\Pr\{\chi^2_{d-2}\ge(d-2)/2\}\ge1/12$.  Hence, with probability at least $1/1152$,
\[
 \lVert C/\alpha-I_m\rVert_{\rm op}
 \ge\sqrt{\frac{(d-2)\gamma_0}{8n\alpha^2}}.
\]
The scaled integrals give
\[
 \frac{\gamma_0}{\alpha^2}\ge
 \frac{m_0(2p)}{2\varphi(0)m_0(p)^2}r,
\]
which proves \eqref{eq:explicitlower}.  If $n\gamma_0<H^2$, then $n=O_p(r)$ and the slab bound
\eqref{eq:lower} supplies a fixed failure probability after decreasing $\delta_p$.

For the confidence term, take $A=W/\alpha$ in Lemma~\ref{lem:gaussianproducttail}.  The finite-radius brackets give
\begin{equation}
 \rho=\frac{\E W^2}{\alpha^2}\ge L_pr,
 \qquad
 \chi=\frac{\E W^4}{(\E W^2)^2}\le K_pr.
 \label{eq:rhochiexplicit}
\end{equation}
Indeed, use $\alpha\le\varphi(0)m_0(p)/r$,
$\E W^2\ge\varphi(0)m_0(2p)/(2r)$, and
$\E W^4\le\varphi(0)m_0(4p)/r$.  If $n<3K_pr$, the slab bound \eqref{eq:lower} gives failure at least
$\delta_p^{\rm hc}$.  Thus success for $\delta<\delta_p^{\rm hc}$ forces $n\ge3K_pr$, and the last sentence of
Lemma~\ref{lem:gaussianproducttail} gives failure at least
\[
 \frac3{80}\exp\!\left(-\frac{2n\varepsilon^2}{L_pr}\right).
\]
This proves \eqref{eq:explicithighconfidence}.  Combining it with the dimension bound by
$\max\{a,b\}\ge(a+b)/2$ proves \eqref{eq:matchinglower} with explicit admissible ranges.
\end{proof}

For reproducibility, choosing $a_p$ by equality in the defining tail condition gives the following numerical
constants.  The tiny $\delta_p^{\rm hc}$ values reflect deliberately crude slab and fourth-moment bounds; they are
proof ranges, not estimates of practical behavior.
\begin{center}
\begin{tabular}{ccccc}
\toprule
$p$ & $a_p$ & $L_p$ & $K_p$ & $\delta_p^{\rm hc}$\\
\midrule
$1$ & $6.255617$ & $0.208886$ & $2.578246$ & $1.4661\times10^{-34}$\\
$2$ & $3.517898$ & $0.322281$ & $3.817398$ & $6.0017\times10^{-29}$\\
\bottomrule
\end{tabular}
\end{center}

The theorem is about the fixed-oracle weighted Gram matrix, not a minimax estimator.  Nor should it be exported
unchanged to weaker targets.  Direct scalar Bernstein gives an $O_p(r/\varepsilon^2)$ radial-entry upper bound at
constant confidence; the empty-slab argument gives only the fixed-accuracy $\Omega_p(r)$ obstruction recorded
above.  The ordinary bottom eigenspace behaves differently: in the iterated extreme-saturation limit it is governed
by sample order statistics rather than relative Gram resolution.

The algebraic flag mechanism below is an instance of classical scaled-SVD asymptotics
\cite{stewart1984}.  The contribution here is its logistic ordering by distance to the transition hyperplane, the
resulting bottom projector, and the exact stochastic angle law and complexity.

\begin{theorem}[spectral lexicography and exact limiting angle law]
\label{thm:spectrallexicography}
Fix $p>0$, $d\ge2$, $n\ge d$, and $u\in S^{d-1}$.  Let $X_i\stackrel{\rm iid}{\sim}N(0,I_d)$,
$Z_i=u^\top X_i$, and
\[
 \widehat H_r=\frac1n\sum_{i=1}^n\sigma'(rZ_i)^pX_iX_i^\top.
\]
Let $\pi$ be the almost surely unique permutation for which
$|Z_{\pi(1)}|<\cdots<|Z_{\pi(n)}|$, and put
$V_k=\operatorname{span}\{X_{\pi(1)},\ldots,X_{\pi(k)}\}$.  Almost surely, for every $k<d$ and all sufficiently
large $r$, a gap separates the top $k$ eigenvalues.  If $E_{r,k}$ denotes that cluster's projector, then
\begin{equation}
 \lVert E_{r,k}-P_{V_k}\rVert_{\rm op}\longrightarrow0,
 \qquad k=1,\ldots,d-1.
 \label{eq:spectralflag}
\end{equation}
For all sufficiently large $r$, the bottom eigenvalue is simple almost surely.  Denoting its rank-one projector by
$\widehat P_r^{\rm bot}$, one has convergence in operator norm,
\begin{equation}
 \lVert\widehat P_r^{\rm bot}-P_\infty\rVert_{\rm op}\longrightarrow0,
 \qquad P_\infty:=P_{V_{d-1}^{\perp}}.
 \label{eq:bottomlimit}
\end{equation}

Write $m=d-1$, let $A_{(j)}=|Z_{\pi(j)}|$, and let
$\zeta=(Z_{\pi(1)},\ldots,Z_{\pi(m)})^\top$.  If $\theta_\infty$ is the acute angle between the range of
$P_\infty$ and $u$, then the following identity is exact:
\begin{equation}
 \tan\theta_\infty\ \stackrel{d}=\ \frac{\lVert\zeta\rVert}{|G|},
 \qquad G\sim N(0,1),\quad G\ \hbox{independent of }\zeta.
 \label{eq:exactanglelaw}
\end{equation}
Let $c_+=\sqrt{2/\pi}$ and $c_-=2\varphi(1)=c_+e^{-1/2}$.  For $m\ge2$,
\begin{equation}
 \Pr\!\left\{\lVert\zeta\rVert\ge
 \frac{m^{3/2}}{6\sqrt2c_+n}\right\}\ge\frac12.
 \label{eq:orderlower}
\end{equation}
If $n\ge4m/c_-$, then
\begin{equation}
 \Pr\!\left\{\lVert\zeta\rVert\le
 \frac{4m^{3/2}}{c_-n}\right\}\ge1-e^{-9m/8}.
 \label{eq:orderupper}
\end{equation}
Let $\theta_r$ be the acute angle between the range of $\widehat P_r^{\rm bot}$ and $u$.  For $d\ge3$ it follows
that, for $0<\varepsilon\le1/2$ and $0<\delta<c_-/2$,
\begin{equation}
 \lim_{r\to\infty}\Pr\{\sin\theta_r\le\varepsilon\}\ge1-\delta
 \quad\Longrightarrow\quad
 n\ge\frac{e^{-1/2}}{24\sqrt2}\,
 \frac{(d-1)^{3/2}}{\varepsilon\delta}.
 \label{eq:anglelowercomplexity}
\end{equation}
Conversely, if $\delta\ge2e^{-9m/8}$, $n\ge4m/c_-$, and
\begin{equation}
 n\ge8e^{1/2}\frac{m^{3/2}}{\varepsilon\delta},
 \label{eq:angleuppercomplexity}
\end{equation}
then the limiting success probability is at least $1-\delta$.  Thus the limiting sample complexity has matching
order $\Theta(m^{3/2}/(\varepsilon\delta))$ in the displayed confidence range, and
$\Theta(d^{3/2}/\varepsilon)$ at fixed confidence.
\end{theorem}

\begin{proof}
Order the samples as $x_j=X_{\pi(j)}$ and the weights as $w_j(r)=\sigma'(rZ_{\pi(j)})^p$.  Since
\[
 \sigma'(s)=\frac{e^{-|s|}}{(1+e^{-|s|})^2},
\]
one has $w_k(r)/w_j(r)\to0$ exponentially whenever $j<k$.  Let $B_r$ have columns
$\sqrt{w_j(r)}x_j$, so $\widehat H_r=n^{-1}B_rB_r^\top$.  Cauchy--Binet on the $k$th exterior power gives
\[
 \bigwedge^k\widehat H_r
 =n^{-k}\!\sum_{|I|=k}\left(\prod_{i\in I}w_i(r)\right)
 (x_{i_1}\wedge\cdots\wedge x_{i_k})(x_{i_1}\wedge\cdots\wedge x_{i_k})^\top.
\]
After division by $\prod_{j=1}^kw_j(r)$, only $I=\{1,\ldots,k\}$ survives.  Its wedge is nonzero almost surely.
The top projector of the exterior power therefore converges to its rank-one span.  Because its top eigenvector is
the wedge of the top $k$ eigenvectors of $\widehat H_r$, this proves \eqref{eq:spectralflag}; taking orthogonal
complements at $k=d-1$ proves \eqref{eq:bottomlimit}.

Rotate so $u=e_d$ and write $X=(Y,Z)$.  Let $Y_*$ be the $m\times m$ matrix whose rows are the tangential
coordinates $Y_{\pi(j)}^\top$, $j\le m$.  Selection depends only on the $Z_i$, so $Y_*$ is an iid standard
Gaussian matrix independent of $\zeta$.  A normal to the selected span is proportional to
$(-Y_*^{-1}\zeta,1)$, hence $\tan\theta_\infty=\lVert Y_*^{-1}\zeta\rVert$.  Conditional on $\zeta$, rotate
row space so that $\zeta$ becomes $\lVert\zeta\rVert e_1$.  If $v$ is a unit normal to rows $2,\ldots,m$ of the
rotated Gaussian matrix and $a_1$ is its first row, then
$Y_*^{-1}\zeta$ has norm $\lVert\zeta\rVert/|a_1^\top v|$.  Conditional on the remaining rows,
$a_1^\top v\sim N(0,1)$, proving \eqref{eq:exactanglelaw}.

It remains to bound the first $m$ half-normal order statistics.  Their distribution function
$F(t)=2\Phi(t)-1$ satisfies
\begin{equation}
 F(t)\le c_+t\quad(t\ge0),
 \qquad
 F(t)\ge c_-t\quad(0\le t\le1).
 \label{eq:halfnormalbounds}
\end{equation}
For $k=\lfloor m/2\rfloor$ and $t=k/(2c_+n)$, the count
$N_t=\#\{i:|Z_i|\le t\}$ has mean at most $k/2$.  Markov's inequality gives
$\Pr\{A_{(k)}>t\}\ge1/2$.  Since $k\ge m/3$ and at least $m/2$ of the first $m$ order statistics exceed
$A_{(k)}$, this proves \eqref{eq:orderlower}.  For the other direction take $t=4m/(c_-n)\le1$.
Now $\E N_t\ge4m$, and a Chernoff bound gives
$\Pr\{N_t<m\}\le e^{-9m/8}$.  On the complement,
$\lVert\zeta\rVert\le\sqrt mA_{(m)}\le4m^{3/2}/(c_-n)$, proving \eqref{eq:orderupper}.

Finally, $\Pr\{|G|\le x\}\ge c_-\min\{x,1\}$ and
$\Pr\{|G|\le x\}\le c_+x$.  Combining the first inequality with \eqref{eq:orderlower} yields, for every $\tau>0$,
\[
 \Pr\{\tan\theta_\infty>\tau\}\ge
 \frac{c_-}{2}\min\!\left\{\frac{m^{3/2}}{6\sqrt2c_+n\tau},1\right\}.
\]
Since $\sin\theta\le\varepsilon$ implies
$\tan\theta\le\varepsilon/\sqrt{1-\varepsilon^2}\le2\varepsilon$, this proves
\eqref{eq:anglelowercomplexity}.  For the upper bound, with probability at least $1-q-e^{-9m/8}$,
\[
 \tan\theta_\infty\le\frac{4c_+}{c_-q}\frac{m^{3/2}}n.
\]
Set $q=\delta/2$ and use $c_+/c_-=e^{1/2}$ to obtain \eqref{eq:angleuppercomplexity}.  Almost-sure projector
convergence and continuity of the limiting angle law justify the displayed probability limit.
\end{proof}

\begin{figure}[H]
 \centering
 \includegraphics[width=.96\linewidth]{../figures/spectral_lexicography.pdf}
 \caption{Fixed-seed diagnostics for Theorem~\ref{thm:spectrallexicography}.  Left: the limiting angular loss
 follows the predicted $n^{-1}$ scale at fixed dimension (the theorem gives matching
 $m^{3/2}/(\varepsilon\delta)$ complexity in its stated confidence range).  Right: finite-$r$ bottom projectors approach the sample-wise
 lexicographic limit; the proof, not this simulation, establishes convergence.}
 \label{fig:spectrallexicography}
\end{figure}

The order of limits in Theorem~\ref{thm:spectrallexicography} is essential: $d,n,p,\varepsilon,\delta$ are fixed
while $r\to\infty$, and only then is the family of limit experiments compared across $n,d$.  The order-statistic
bounds are uniform under their stated conditions, but \eqref{eq:bottomlimit} has no rate uniform in a joint sequence
$n=n(r)$ or $d=d(r)$.  Such a theorem would need simultaneous control of adjacent half-normal spacings and the
conditioning of $Y_*$.  In particular, an $\Omega(rd)$ angular necessity cannot hold uniformly through extreme
saturation; it remains the scale of the sufficient relative-Gram event, not of the weaker limiting projector target.

\subsection{Uniform all-power sensitivity on a full dyadic shell}

The pointwise theorem admits a proof-complete full-dyadic-shell extension at the price of an elementary covering
logarithm.
For $G_p(w)=\E[h_p(w^\top X)XX^\top]$ and its empirical analogue $\widehat G_{p,n}(w)$, set
\[
 B_{n,\delta}=\sqrt d+\sqrt{2\log(2n/\delta)},\quad
 \mu_{3,d}=[d(d+2)]^{3/4},\quad L_p=p4^{-p},
\]
\[
 b_p=\frac{\varphi(0)m_2(p)}{16},\quad
 A_{p,n,\delta}=L_p(B_{n,\delta}^3+\mu_{3,d}),\quad
 \eta=\min\left\{R,\frac{\varepsilon b_p}{32A_{p,n,\delta}R^3}\right\},\quad
 N=\left(1+\frac{4R}{\eta}\right)^d.
\]

\begin{theorem}[uniform relative sensitivity on a full dyadic shell]
\label{thm:uniformshell}
Fix $p>0$, $R\ge R_p$, $0<\varepsilon\le1/2$, and $0<\delta<1/2$.  If
\begin{equation}
 n\ge128C_p\frac{R}{\varepsilon^2}
 \left[d+\log\left(\frac{24N}{\delta}\right)\right],
 \label{eq:uniformn}
\end{equation}
then, with probability at least $1-\delta$, simultaneously for every $w$ with
$R\le\lVert w\rVert\le2R$,
\begin{equation}
 (1-\varepsilon)G_p(w)\preceq\widehat G_{p,n}(w)
 \preceq(1+\varepsilon)G_p(w).
 \label{eq:uniformgram}
\end{equation}
The computable condition is implicit in $n$ only through $B_{n,\delta}$; no optimality is claimed for its covering
logarithm.
\end{theorem}

\begin{proof}
Since $\lVert h_p'\rVert_\infty\le L_p$,
\[
 \lVert h_p(w^\top x)xx^\top-h_p(v^\top x)xx^\top\rVert_{\rm op}
 \le L_p\lVert w-v\rVert\lVert x\rVert^3.
\]
With probability at least $1-\delta/2$, all sampled Gaussian norms are at most $B_{n,\delta}$.  On this event
$D_n=\widehat G_{p,n}-G_p$ is $A_{p,n,\delta}$-Lipschitz in operator norm.  Cover the shell by an $\eta$-net of
size at most $N$ and apply Corollary~\ref{cor:effective} at every net point with error $\varepsilon/8$ and failure
$\delta/(2N)$; \eqref{eq:uniformn} is exactly a conservative sufficient condition.  On the shell,
$\lambda_{\min}G_p(w)\ge b_pR^{-3}$.  Moving from a net point to $w$ therefore changes both $D_n$ and $G_p$ by
at most $(\varepsilon/32)\lambda_{\min}G_p(w)$.  For every $x$,
\[
 |x^\top D_n(w)x|
 \le\left[\frac\varepsilon8(1+\varepsilon/32)+\frac\varepsilon{32}\right]x^\top G_p(w)x
 \le\varepsilon x^\top G_p(w)x,
\]
which is \eqref{eq:uniformgram}.  The two failure probabilities total at most $\delta$.
\end{proof}

For $p=1$, $\widehat G_{1,n}(w)$ is exactly the logistic-likelihood Hessian and is label-free.  For $p\ne1$ it is
a sensitivity Gram, not that likelihood Hessian.  Earlier work proves uniform local constant-factor sandwiches and
an optimal-scale one-sided $p=1$ lower bound \cite{ostrovskii2021,chardon2024}; the narrower addition here is a
full dyadic shell, all $p>0$, and tunable bilateral precision.  It is not by itself an estimator theorem.

\begin{figure}[H]
 \centering
 \includegraphics[width=\textwidth]{../figures/finite_sample_resolution.pdf}
 \caption{Fixed-seed diagnostic at $d=12$, $r=6$, with 320 repetitions per sample size.  Solid curves are medians
 and dashed curves are 90th percentiles.  The first three panels display tangential relative error, ordinary
 smallest-eigenvalue relative error, and radial eigenspace angle.  The fourth compares the exact empty-slab
 probability with its simulated frequency; exact values below $10^{-8}$ are clipped only in the plot and remain
 untruncated in the JSON certificate.  The normalization $n/(rd)$ is the sharp scale for full bilateral
 relative Loewner loss; the ordinary eigenvalue and angle panels are diagnostics and make no claim that this is
 their optimal threshold.}
 \label{fig:finite}
\end{figure}

\section{What the law costs}

\subsection{Stationary fixed-step gradient descent}

For either logistic profile, assume $d\ge2$ and that one constant scalar step size $\eta$ is reused at every
iteration.  At the teacher, the
Jacobian of one gradient-descent step is $I-\eta H_h$.  Its radial eigenvalue is
$1-\eta\beta$ and its tangential eigenvalue is $1-\eta\alpha$.  The best scalar step for this local quadratic model is
\begin{equation}
 \eta_\star=\frac{2}{\alpha+\beta},\qquad
 \rho_\star=\frac{\alpha-\beta}{\alpha+\beta}=\frac{\kappa-1}{\kappa+1}.
 \label{eq:optrate}
\end{equation}
Consequently $-\log\rho_\star\sim2/\kappa$.  For worst-case local error with components allowed in both
eigenspaces, the e-folding iteration count
$T_{\rm e}:=(-\log\rho_\star)^{-1}$ therefore has the exact leading scaling
\begin{equation}
 T_{{\rm e},\rm sq}\sim\frac{3}{2(\pi^2-6)}r^2,\qquad
 T_{{\rm e},\rm B}\sim\frac3{2\pi^2}r^2
 \label{eq:iterationtax}
\end{equation}
for the optimally tuned stationary scalar-step linearization.  This is not a global iteration bound or an
algorithm-independent obstruction: away from the teacher the Hessian contains residual terms, and nonstationary
polynomial methods can exploit the two-point spectrum.  Indeed, the scalar schedule
$\eta_1=1/\alpha$, $\eta_2=1/\beta$ makes
$(I-\eta_2H)(I-\eta_1H)=0$ in this exact local quadratic.  Natural-gradient or matrix-preconditioned updates can
also remove the local condition number by acting differently on the two eigenspaces.

\subsection{Estimation}

For $n$ independent Gaussian-output observations with known variance $\tau^2$, the Fisher information is
$nH_{h_2}/\tau^2$.  The Cram\'er--Rao inequality gives, along any unit tangential vector $v\perp u$ and the radial
vector $u$,
\begin{equation}
 \operatorname{Var}(v^\top\widehat w)\ge\frac{\tau^2}{n\alpha_2(r)},\qquad
 \operatorname{Var}(u^\top\widehat w)\ge\frac{\tau^2}{n\beta_2(r)}.
 \label{eq:crgauss}
\end{equation}
The ratio of the displayed radial lower bound to the tangential lower bound is $\kappa_2(r)$.  Hence tangential
variance is bounded below at order $r/n$, while radial variance is bounded below at order
$r^3/n$.  The analogous Bernoulli bounds replace $h_2$ by $h_1$.  The same quadratic ratio therefore appears in
these coordinatewise information lower bounds and in stationary fixed-step descent, although the latter can be
removed by a nonstationary schedule.

\section{Bias and universality}

A conventional neuron includes a bias $b$.  At $b=0$ the augmented parameter $(w,b)$ has Fisher/Hessian block
\[
 \E\!\left[h(rZ)
 \begin{pmatrix}XX^\top&X\\X^\top&1\end{pmatrix}\right].
\]
For even $h$, the radial--bias cross term $\E[Zh(rZ)]$ vanishes.  The bias eigenvalue equals $\alpha_h(r)$, so
the spectrum contains one radial eigenvalue $\beta_h$ and $d$ tangential/bias eigenvalues $\alpha_h$.  A nonzero
bias preserves the $(d-1)$ tangential eigenspace but replaces the radial and bias coordinates by the explicit
$2\times2$ block
\begin{equation}
 \begin{pmatrix}
  \E[Z^2h(rZ+b)] & \E[Zh(rZ+b)]\\
  \E[Zh(rZ+b)] & \E[h(rZ+b)]
 \end{pmatrix}.
 \label{eq:biasblock}
\end{equation}
The centered theorem is therefore not an artifact of omitting the bias; it is the diagonal point of a fully explicit
two-coordinate extension.

The exponent two is broader than the logistic curve.  Theorem~\ref{thm:universal} shows that any localized
sensitivity profile with finite positive $m_0$ and $m_2$ has the same $r^2$ anisotropy, with only the constant
$m_0/m_2$ depending on the activation and observation model.  Profiles with heavy sensitivity tails or vanishing
second moment fall outside this universality class and can have different laws.



\section{Relation to prior work}

Table~\ref{tab:priority} states the claim boundary before the detailed discussion.
\begin{table}[H]
\centering
\small
\begin{tabular}{>{\raggedright\arraybackslash}p{0.27\textwidth}
                >{\raggedright\arraybackslash}p{0.29\textwidth}
                >{\raggedright\arraybackslash}p{0.34\textwidth}}
\toprule
Prior result & Result here & Exact difference \\
\midrule
Fixed-order L\"owner/Schwarzian criterion and serial closure
\cite{kozlovski2008,cook2023}
& closed logistic determinant and certified $2\times2$ witness
& explicit specialization and certificate; not a new criterion or closure law \\[3pt]
One-unit Gaussian Fisher split for $\phi'^2$ \cite{amari2019}
& spherical local bridge and Gaussian all-$p$ laws
& coefficient cross-identification and extensions; not a new one-unit decomposition \\[3pt]
Bernoulli exponents and finite-sample logistic Hessian bounds
\cite{chen2026gd,ostrovskii2021,chardon2024,fisher2023}
 & pointwise and shell-uniform relative Gram bounds, plus matching pointwise complexity
 & tunable precision and all powers; not the first bilateral, uniform, or one-sided Hessian bound \\[3pt]
 Scaled-SVD flags \cite{stewart1984} and Gaussian logistic direction/Hessian theory
 \cite{hsu2024,chardon2024,chen2026gd,chen2026minimax}
 & spectral lexicography and exact extreme-saturation angle law
 & flag mechanism is classical; logistic sample ordering, angle law and complexity are the refinement; not labeled
 direction estimation or a joint finite-$r$ rate \\[3pt]
 Classical spherical representation \cite{cambanis1981} and regular-design reference exponents \cite{chardon2024}
 & exact $Q_R$ constants and three inverse-radius phases
 & leading constants and critical logarithm; not a new spherical representation or first
non-Gaussian exponent \\[3pt]
Generalized-logistic transforms and cumulants \cite{ojo2004,lee2025}
& Gamma/polygamma moments inside neuronal anisotropy
& transform is classical; the unified anisotropy, endpoint laws and monotonicity are the contribution \\[3pt]
Lam's historical v2 Fisher--Schwarzian construction \cite{lam2026v2}
& L\"owner--anisotropy coefficient identity
& different objects and coefficient; no claim to the first general Fisher--Schwarzian connection \\
\bottomrule
\end{tabular}
\caption{Result-by-result priority boundary.}
\label{tab:priority}
\end{table}

L\"owner matrices characterize fixed-order matrix monotonicity; modern treatments include Hiai--Sano and
Hein\"avaara \cite{hiai2010,heinavaara2019}.  Kozlovski and Sands prove the relevant classical equivalence between
$Sf\ge0$ and order-two matrix monotonicity and record the Schwarzian composition law
\cite{kozlovski2008}.  Hence the qualitative logistic obstruction and its serial-chain extension are classical
consequences; negative Schwarzian also implies the all-pair determinant sign through strict convexity of
$(f')^{-1/2}$.  Cook, Hammerlindl and Tucker develop the algebraically equivalent nowhere-coexpanding inequality,
prove serial closure, and explicitly discuss sigmoid activations and one-dimensional neural compositions
\cite{cook2023}.  Neural papers have proposed analytic
quantum activations and trainable matrix activations \cite{maronese2022,liu2021}, but their objectives are circuit
implementation or matrix-valued parameterization, not preservation of PSD order by the spectral logistic function.
In the primary sources we located, we did not find the closed logistic identity \eqref{eq:loewnerdet}, the explicit
certified witness, or the coefficient identity \eqref{eq:bridge} linking the L\"owner defect to Gaussian
radial--tangential anisotropy.  This is a bounded novelty statement, not a claim that no earlier source can exist.

Learning one neuron is already a nontrivial optimization problem.  Yehudai and Shamir study when gradient methods
learn a single neuron under broad activations and input distributions, including positive and negative results
\cite{yehudai2020}.  Diakonikolas et al. treat monotone neurons with adversarial label noise and obtain polynomial
learners for logistic activations under log-concave inputs \cite{diakonikolas2022}.  Wu extends learnability results to
nonmonotone activations \cite{wu2022}, and Vardi et al. show that adding a bias qualitatively changes the ReLU
landscape \cite{vardi2021}.  Recent single-index work analyzes SGD under anisotropic Gaussian inputs
\cite{braun2025}.  Under Gaussian logistic regression, Hsu and Mazumdar study the signal-dependent difficulty of
direction and temperature estimation \cite{hsu2024}.  Chen and Mazumdar identify the Bernoulli radial and
tangential Hessian functions, prove their $r^{-3}$ and $r^{-1}$ orders and show tangential curvature exceeds radial
curvature for every nonzero signal \cite{chen2026gd}; they subsequently prove a minimax norm-estimation rate of
order $\sqrt{r^3/n}$ \cite{chen2026minimax}.  Ostrovskii and Bach prove a uniform bilateral constant-factor
empirical-Hessian sandwich for Gaussian logistic regression \cite[Eq.~(92)]{ostrovskii2021}, while Fisher et al.
give a generic fixed-point standardized-Hessian concentration framework \cite{fisher2023}.  Chardon, Lerasle and
Mourtada prove, for $p=1$, a uniform one-sided empirical Hessian lower bound under the sufficient condition
$n\gtrsim r(d+t)$ \cite[Theorem 6]{chardon2024}.  We therefore do not claim the Bernoulli exponents, the $p=1$
ordering, the effective transition-slab scale, the first bilateral Hessian concentration, or the first finite-sample
lower bound.  Our empirical contribution is the explicit pointwise $(1\pm\varepsilon)$ approximation
\eqref{eq:relativegram} for all $p>0$, its matching bilateral lower complexity, the full-dyadic-shell extension
\eqref{eq:uniformgram}, and the distinct extreme-saturation projector law
\eqref{eq:exactanglelaw}.  The shell theorem pays a conservative covering logarithm and is not asserted to improve the
optimal local one-sided rate; the projector theorem is iterated-limit rather than uniform finite-$r$.  Stewart's
scaled-SVD analysis already supplies the underlying deterministic flag mechanism \cite{stewart1984}.  In the primary
sources searched through August 30, 2026, we did not locate the logistic nearest-hyperplane projector, its exact
angle law, or the resulting order-statistic complexity.  This is a bounded search statement, not a guarantee of
priority.  Recent work already constructs minimax labeled estimators, including a split-sample
direction/norm scheme with bias correction \cite{chen2026minimax}; our oracle Gram theorem is not an estimator.

Fisher geometry and natural gradient are classical \cite{amari1998}.  Karakida et al. study broad Fisher spectral
statistics for random deep networks \cite{karakida2019}; Amari, Karakida and Oizumi also give the direct one-unit
Gaussian $\phi'^2$ decomposition and bias coupling used here \cite{amari2019}.  Our radial contribution is
therefore not the two-eigenspace reduction.  It is the profile-generic moment theorem, nonasymptotic brackets,
closed all-power logistic constants, strict global monotonicity, and their connection through \eqref{eq:bridge} to
the L\"owner order defect.  The generalized-logistic Gamma transform and its polygamma cumulants in
Lemma~\ref{lem:moments} are classical \cite{ojo2004,lee2025}; the claim is their neural-geometric use, ratio
monotonicity and endpoint package, not discovery of the distributional identity.
Lam's v2 preprint connected Schwarzian curvature with Fisher--Rao geometry on manifolds of densities
\cite{lam2026v2}.  The current v3 explicitly corrects and supersedes v1--v2 and removes that architecture
\cite{lam2026v3}.  Neither version contains the one-neuron radial--tangential coefficient or its equality with the
two-point L\"owner defect.  We cite v2 only as a dated historical adjacency; its $L^p$ parameter is unrelated to
the sensitivity exponent in $h_p=\sigma'^p$.

\section{Reproducibility and hostile checks}

The documented sequence in \texttt{repro/README.md} runs the original and v3 checks together with three new
scripts:
\begin{center}
\texttt{radial\_phase\_transition.py},\quad
\texttt{spectral\_lexicography.py},\quad
\texttt{verify\_v4\_additions.py}.
\end{center}
Together they perform the following independent checks:
\begin{enumerate}[leftmargin=2em]
 \item it spot-checks \eqref{eq:loewnerdet} at separated scales and constructs Proposition~\ref{prop:witness} with
 256-bit directed Arb balls;
 \item it checks the spherical fourth-moment coefficient for Gaussian and fixed-radius spherical inputs;
 \item it compares numerical quadrature of $m_0,m_2,m_4$ with the Gamma/polygamma formula for several
 noninteger and integer $p$;
 \item it checks the small-$r$ jet through $r^6$ and the refined large-$r$ expansion;
 \item it evaluates $\alpha,\beta$ directly over a logarithmic radius grid and checks positivity and grid
 monotonicity for $p=1,2$;
 \item it checks \eqref{eq:quantkappa}--\eqref{eq:directfinitebridge} at declared finite scales;
 \item for three declared fixed random seeds it evaluates the empirical radial, tangential and cross blocks and
 verifies that those realized errors lie inside \eqref{eq:blockevent};
 \item it checks both sides of \eqref{eq:alphabounds}--\eqref{eq:betabounds} at declared finite radii;
 \item a second replay evaluates the explicit $C_1,C_2$ majorants and lower-bound constants; and
 \item a fixed-seed replay regenerates Figure~\ref{fig:finite} and its JSON certificate, while the main replay
 regenerates Figure~\ref{fig:saturation} and its certificate;
 \item deterministic adaptive quadrature regenerates Figure~\ref{fig:radialphases} and checks all three exact
 phase constants for isotropic Gamma radii; and
 \item a fixed-seed eigenspace replay regenerates Figure~\ref{fig:spectrallexicography}, while a separate check
 validates the explicit $L_p,K_p$ confidence constants and the sample-wise normal-vector identity.
\end{enumerate}
The numerical monotonicity and finite-radius checks are performed on declared grids; the universal statements are
proved analytically above.  The replay uses Python, python-flint, NumPy, SciPy, and Matplotlib.  It is diagnostic
rather than a substitute for
the analytic proofs, which reduce to the L\"owner determinant identity, the Schwarzian chain rule, Gaussian
orthogonal decomposition, Bernstein/net bounds, an Abelian kernel lemma, exterior powers, elementary order
statistics, an explicit Gaussian lower tail, and closed Fourier moments.
The stable record, versioned source archives, audit reports and manifests are public at
\url{https://arr-research.github.io/papers/ARR-2026-53CTRKDSP685PT51/}.

\section{Limitations and next theorem}

The L\"owner theorem concerns spectral matrix functions, not entrywise vector activations; it does not say that a
standard feedforward neuron reverses coordinatewise order.  The empirical results assume Gaussian input.  The
local Schwarzian coefficient needs fourth moments, whereas the spherical asymptotic law needs transition-layer
regularity; $\E R^{-1}<\infty$ is needed precisely for its quadratic phase, with the other phases covered by
Theorem~\ref{thm:radialphase}.  Covariance isotropy alone is insufficient by the Rademacher example following it.
Elliptical laws can be whitened, but radial/tangential directions then use the covariance metric.

The sharp Gram complexity statement is pointwise for a fixed teacher and bilateral relative Loewner loss.  The
bottom-eigenspace result solves a different, iterated limit: it does not provide a finite-$r$ rate uniform along
$n=n(r)$ or $d=d(r)$.  Closing that joint regime is the strongest next spectral problem.  The dyadic-shell theorem
is uniform but its Euclidean-net logarithm is not quantitatively sharp.  Already in $d=2$ the tangential forms
resemble a marked angular kernel process of bandwidth $1/R$; classical sup-norm kernel laws therefore suggest that
a $\sqrt{\log R}$ penalty is intrinsic \cite{bickel1973,gine2004}.  We do not prove the corresponding marked-process
lower bound.  A model-specific matching shell lower bound and a chaining upper bound remain open.
None of these oracle results is an estimator guarantee using labels or a minimax parameter lower bound.  The
optimization result remains local, and the Cram\'er--Rao statement concerns regular unbiased estimation.

\section*{AI assistance disclosure}
AI tools assisted with public-literature and repository triage, proof stress testing, numerical replay, drafting, and
PDF production.  The author remains responsible for the mathematical claims and any remaining errors.

\begin{thebibliography}{99}
\let\oldbibitem\bibitem
\renewcommand{\bibitem}{\par\oldbibitem}
\setlength{\itemsep}{1pt}
\small
\bibitem{hiai2010}
F.~Hiai and T.~Sano,
\emph{L\"owner matrices of matrix convex and monotone functions},
Journal of the Mathematical Society of Japan 64(2) (2012), 343--364.
\url{https://arxiv.org/abs/1007.2478}

\bibitem{heinavaara2019}
O.~Hein\"avaara,
\emph{Characterizing matrix monotonicity of fixed order on general sets},
preprint (2019).
\url{https://arxiv.org/abs/1906.06155}

\bibitem{kozlovski2008}
O.~Kozlovski and D.~Sands,
\emph{Higher order Schwarzian derivatives in interval dynamics},
Fundamenta Mathematicae 206 (2009), 217--239.
\url{https://arxiv.org/abs/0812.2646}

\bibitem{cook2023}
A.~Cook, A.~Hammerlindl and W.~Tucker,
\emph{Nowhere coexpanding functions},
Chaos 33 (2023), 123105.
\url{https://arxiv.org/abs/2303.12814}

\bibitem{maronese2022}
M.~Maronese, C.~Destri and E.~Prati,
\emph{Quantum activation functions for quantum neural networks},
Quantum Information Processing 21 (2022), 128.
\url{https://arxiv.org/abs/2201.03700}

\bibitem{liu2021}
Z.~Liu, S.~Cao, Y.~Li and L.~Zikatanov,
\emph{Neural networks with trainable matrix activation functions},
Journal of Machine Learning for Modeling and Computing 6(2) (2025), 1--11.
\url{https://arxiv.org/abs/2109.09948}

\bibitem{cambanis1981}
S.~Cambanis, S.~Huang and G.~Simons,
\emph{On the theory of elliptically contoured distributions},
Journal of Multivariate Analysis 11(3) (1981), 368--385.
\url{https://doi.org/10.1016/0047-259X(81)90082-8}

\bibitem{stewart1984}
G.~W.~Stewart,
\emph{On the asymptotic behavior of scaled singular value and QR decompositions},
Mathematics of Computation 43(168) (1984), 483--489.
\url{https://doi.org/10.1090/S0025-5718-1984-0758196-7}

\bibitem{yehudai2020}
G.~Yehudai and O.~Shamir,
\emph{Learning a Single Neuron with Gradient Methods},
Proceedings of Machine Learning Research 125 (2020), 3756--3786.
\url{https://proceedings.mlr.press/v125/yehudai20a.html}

\bibitem{diakonikolas2022}
I.~Diakonikolas, V.~Kontonis, C.~Tzamos and N.~Zarifis,
\emph{Learning a Single Neuron with Adversarial Label Noise via Gradient Descent},
Proceedings of Machine Learning Research 178 (2022), 4313--4361.
\url{https://proceedings.mlr.press/v178/diakonikolas22c.html}

\bibitem{wu2022}
L.~Wu,
\emph{Learning a Single Neuron for Non-monotonic Activation Functions},
Proceedings of Machine Learning Research 151 (2022), 4178--4197.
\url{https://proceedings.mlr.press/v151/wu22c.html}

\bibitem{vardi2021}
G.~Vardi, G.~Yehudai and O.~Shamir,
\emph{Learning a Single Neuron with Bias Using Gradient Descent},
Advances in Neural Information Processing Systems 34 (2021), 28690--28700.
\url{https://arxiv.org/abs/2106.01101}

\bibitem{braun2025}
G.~Braun, M.~H.~Quang and M.~Imaizumi,
\emph{Learning a Single Index Model from Anisotropic Data with Vanilla Stochastic Gradient Descent},
Proceedings of Machine Learning Research 258 (2025), 1216--1224.
\url{https://proceedings.mlr.press/v258/braun25a.html}

\bibitem{hsu2024}
D.~Hsu and A.~Mazumdar,
\emph{On the Sample Complexity of Parameter Estimation in Logistic Regression with Normal Design},
Proceedings of Machine Learning Research 247 (2024), 2418--2437.
\url{https://proceedings.mlr.press/v247/hsu24a.html}

\bibitem{chen2026gd}
J.~Chen and A.~Mazumdar,
\emph{Finite-Sample Performance of Gradient Descent in Logistic Regression with Gaussian Design},
arXiv:2606.21683 (2026).
\url{https://arxiv.org/abs/2606.21683}

\bibitem{chen2026minimax}
J.~Chen and A.~Mazumdar,
\emph{Minimax Optimal Estimator and Improved Error Rate for the MLE in Logistic Regression with Gaussian Design},
arXiv:2608.17260 (2026).
\url{https://arxiv.org/abs/2608.17260}

\bibitem{chardon2024}
H.~Chardon, M.~Lerasle and J.~Mourtada,
\emph{Finite-sample performance of the maximum likelihood estimator in logistic regression},
arXiv:2411.02137v3 (first posted 2024; v3 2026).
\url{https://arxiv.org/abs/2411.02137v3}

\bibitem{ostrovskii2021}
D.~M.~Ostrovskii and F.~Bach,
\emph{Finite-sample analysis of M-estimators using self-concordance},
Electronic Journal of Statistics 15(1) (2021), 326--391.
\url{https://doi.org/10.1214/20-EJS1780}

\bibitem{bickel1973}
P.~J.~Bickel and M.~Rosenblatt,
\emph{On some global measures of the deviations of density function estimates},
Annals of Statistics 1(6) (1973), 1071--1095.
\url{https://doi.org/10.1214/aos/1176342558}

\bibitem{gine2004}
E.~Gin\'e, V.~Koltchinskii and L.~Sakhanenko,
\emph{Kernel density estimators: convergence in distribution for weighted sup-norms},
Probability Theory and Related Fields 130 (2004), 167--198.
\url{https://doi.org/10.1007/s00440-003-0314-z}

\bibitem{fisher2023}
J.~Fisher, L.~Liu, K.~Pillutla, Y.~Choi and Z.~Harchaoui,
\emph{Influence Diagnostics under Self-concordance},
Proceedings of Machine Learning Research 206 (2023), 10028--10076.
\url{https://proceedings.mlr.press/v206/fisher23a.html}

\bibitem{amari1998}
S.-i.~Amari,
\emph{Natural Gradient Works Efficiently in Learning},
Neural Computation 10 (1998), 251--276.
\url{https://doi.org/10.1162/089976698300017746}

\bibitem{karakida2019}
R.~Karakida, S.~Akaho and S.-i.~Amari,
\emph{Universal Statistics of Fisher Information in Deep Neural Networks: Mean Field Approach},
Proceedings of Machine Learning Research 89 (2019), 1032--1041.
\url{https://proceedings.mlr.press/v89/karakida19a.html}

\bibitem{amari2019}
S.-i.~Amari, R.~Karakida and M.~Oizumi,
\emph{Fisher Information and Natural Gradient Learning in Random Deep Networks},
Proceedings of Machine Learning Research 89 (2019), 694--702.
\url{https://proceedings.mlr.press/v89/amari19a.html}

\bibitem{lee2025}
C.~J.~Lee, A.~Zito, H.~Sang and D.~B.~Dunson,
\emph{Logistic-beta processes for dependent random probabilities with beta marginals},
Bayesian Analysis 20(4) (2025), 1345--1369.
\url{https://doi.org/10.1214/25-BA1541}

\bibitem{ojo2004}
M.~O.~Ojo and A.~K.~Olapade,
\emph{On a Six-Parameter Generalized Logistic Distribution},
Kragujevac Journal of Mathematics 26 (2004), 31--38.
\url{https://imi.pmf.kg.ac.rs/kjm/pub/12616736649184_5.pdf}

\bibitem{lam2026v2}
H.~P.~G.~Lam,
\emph{Real Bers embedding on the line: Fisher--Rao linearization, Schwarzian curvature, and scattering coordinates},
arXiv:2602.07373v2 (2026).
\url{https://arxiv.org/abs/2602.07373v2}

\bibitem{lam2026v3}
H.~P.~G.~Lam,
\emph{Zero-energy scattering and the real Bers image on the line},
arXiv:2602.07373v3 (2026), correcting and superseding v1--v2.
\url{https://arxiv.org/abs/2602.07373v3}
\end{thebibliography}

\end{document}
